% \iffalse meta-comment
%
% Copyright (C) 2017- by Yanshuo Chu <yanshuoc@gmail.com>
%
% This file may be distributed and/or modified under the
% conditions of the LaTeX Project Public License, either version 1.3a
% of this license or (at your option) any later version.
% The latest version of this license is in:
%
% http://www.latex-project.org/lppl.txt
%
% and version 1.3a or later is part of all distributions of LaTeX
% version 2004/10/01 or later.
%
% \fi
%
% ^^A ======================================================================
% ^^A hithesis 的 book 家族实现
% ^^A
% ^^A 毕业论文类 hithesisbook 的全部代码。
% ^^A
% ^^A 类文件头（\NeedsTeXFormat、\ProvidesExplClass 等）留在 hithesis.dtx，
% ^^A hithesis.ins 里每个 \file 先从那里取头，再从本文件取实现。
% ^^A ======================================================================
%
% ^^A ======================================================================
% ^^A Module thesiscls: hit-thesis 类主体（选项透传 + LoadClass + 模块加载 + AtEndOfClass）
% ^^A ======================================================================
% ^^A ======================================================================
% ^^A Module thesis-options: 文档类选项的声明与后处理（type/campus/fontset 等）（hit-thesis）
% ^^A ======================================================================
%    \begin{macrocode}
%<*thesis-options>
%    \end{macrocode}
%
% 此处添加book控制项，用于cfg中定理等控制
% \changes{v3.0.0}{2020/05/21}{此处添加book控制项，用于cfg中定理等控制}
%    \begin{macrocode}
%    \end{macrocode}
%
% 键名带连字符时（如 \option{cjk-font}）不能直接拿来拼控制序列，另给一个把键名与
% 变量名分开的版本。
%    \begin{macrocode}

\hit@declare@choice@option{type}{}{bachelor,master,doctor,postdoc}{}
%    \end{macrocode}
%
% \changes{v3.1d}{2025/03/03}{模板默认是 \option{bachelor}，为页脚做准备。}
%    \begin{macrocode}
\bool_lazy_any:nF
  {
    { \bool_if_p:N \g__hit_master_bool }
    { \bool_if_p:N \g__hit_doctor_bool }
    { \bool_if_p:N \g__hit_postdoc_bool }
  }
  { \bool_gset_true:N \g__hit_bachelor_bool }
%    \end{macrocode}
%
% 此处设置校区，没有明确给出威海或者深圳时候，默认为哈尔滨校区。
% \changes{v2.0.0}{2018/06/14}{此处添加geometry选项}
% \changes{v3.0.0}{2020/05/20}{此处添加博后选项}
%    \begin{macrocode}
\hit@declare@choice@option{campus}{}{shenzhen,weihai,harbin}{}
\bool_lazy_or:nnF
  { \bool_if_p:N \g__hit_weihai_bool } { \bool_if_p:N \g__hit_shenzhen_bool }
  { \bool_gset_true:N \g__hit_harbin_bool }
%    \end{macrocode}
%
% 此处添加语言选项
% \changes{v3.0.0}{2020/05/20}{此处添加语言选项}
%    \begin{macrocode}
\hit@declare@choice@option{lang}{}{en,zh}{}
\bool_if:NF \g__hit_en_bool { \bool_gset_true:N \g__hit_zh_bool }
%    \end{macrocode}
%
% 设置版芯，由于窝工版芯歧义。
% \changes{v2.0.0}{2018/06/14}{此处添加geometry选项}
%    \begin{macrocode}
\hit@declare@choice@option{newgeometry}{newgeometry_}{one,two,no}{}
%    \end{macrocode}
%
% 设置字距选项，因为word和规范不一致
% \changes{v3.0.13}{2020/11/10}{设置字距选项，因为word和规范不一致}
%    \begin{macrocode}
\hit@declare@choice@option{zijv}{ziju_}{regu,word}{}
%    \end{macrocode}
%
% \changes{v3.1c}{2023/04/16}{设置中文参考文献中相邻两个标号使用逗号“,”还是短线“-”。\issue[137]}
%    \begin{macrocode}
\hit@declare@choice@option{citetwo}{citetwo_}{comma,endash}{endash}
%    \end{macrocode}
%
% \changes{v3.2a}{0000/00/00}{设置参考文献列表中的作者是否全大写和最大显示数量}
% \changes{v3.2a}{0000/00/00}{设置参考文献列表标号的对齐方式}
%    \begin{macrocode}

\bool_new:N \g_bib_authorupper_bool
\bool_set_false:N \g_bib_authorupper_bool
\keys_define:nn { hit / class }
  {
    bibauthorupper .code:n =
      {
        \str_if_eq:nnT {#1} {full}
          { \bool_set_true:N \g_bib_authorupper_bool }
      } ,
  }

\int_new:N \g_bib_maxauthor_int
\int_set:Nn \g_bib_maxauthor_int {3}
\keys_define:nn { hit / class }
  {
    bibmaxauthor .code:n =
      {
        \int_set:Nn \g_bib_maxauthor_int {#1}
        \int_if_zero:nT { \g_bib_maxauthor_int }
          { \int_gdecr:N \g_bib_maxauthor_int }
      } ,
  }

\str_new:N \g_bib_labelalign_str
\keys_define:nn { hit / class }
  { biblabelalign .code:n = { \str_set:Nn \g_bib_labelalign_str {#1} } }

%    \end{macrocode}
%
% 目录中英文是否用 Arial 字体（默认关闭）。
%    \begin{macrocode}
\hit@declare@boolean@option{arialtoc}{false}
%    \end{macrocode}
%
% 章节标题中的英文是否用 Arial 字体（默认打开）。
%    \begin{macrocode}
\hit@declare@boolean@option{arialtitle}{false}
%    \end{macrocode}
%
% 中文本科模板中章标题是否加粗，默认是
% \changes{v3.0.01}{2020/06/06}{中文本科模板中章标题是否加粗，默认是}
%    \begin{macrocode}
\hit@declare@boolean@option{chapterbold}{true}
%    \end{macrocode}
%
% \changes{v1.0.03}{2017/08/29}{默认开启raggedbottom}
% \option{raggedbottom} 选项（默认开启）。如果不开启这个选项，会出现一页中尽量上
% 下对齐，段的间距大。如果开启，尽量使段间距保持一致，页面底部出现空白。
%    \begin{macrocode}
\hit@declare@boolean@option{raggedbottom}{true}
%    \end{macrocode}
%
%
% 字体间距设置（默认关闭）。
%    \begin{macrocode}
\hit@declare@boolean@option{glue}{false}
%    \end{macrocode}
%
% 文科生四级目录设置（默认关闭）。
%    \begin{macrocode}
\hit@declare@boolean@option{tocfour}{false}
%    \end{macrocode}
%
% \changes{v3.2a}{0000/00/00}{目录三个选项的旧名改写同一批三态变量，保证新旧
%   两套名字结果一致。变量要在 \cs{ProcessKeyOptions} 之前声明，因为这几个
%   键也能作类选项。}
% 三个三态变量在这里声明，旧名的处理器直接写它们。新名在 \texttt{hit/toc} 族里
% 注册，写的是同一批变量，所以两套名字可以混用、后设的覆盖先设的。
%    \begin{macrocode}

%    \end{macrocode}
%
% 目录中“目录”位置是否空行（默认开启）。
%    \begin{macrocode}
\hit@declare@boolean@option{tocblank}{true}
%    \end{macrocode}
%
% 章标题是否悬挂居中（默认开启）
%    \begin{macrocode}
\hit@declare@boolean@option{chapterhang}{true}
%    \end{macrocode}
%
% 是否是全日制学生（默认是）。
%    \begin{macrocode}
\hit@declare@boolean@option{fulltime}{true}
%    \end{macrocode}
%
% \changes{v3.2a}{0000/00/00}{增加选项：英文内封标题是否根据长度自动调整字号}
% 英文内封标题是否根据长度自动调整字号（默认关闭，即使用二号字）。
%    \begin{macrocode}
\hit@declare@boolean@option{etitleauto}{false}
%    \end{macrocode}
%
% \changes{v3.2a}{0000/00/00}{\option{subtitle} 从布尔改成三态，默认 \texttt{auto}：
%   \option{title} 里写了方括号就排副标题，不必再单独打开开关}
% 要不要排副标题。默认 \texttt{auto}，即 \option{title} 的取值里带了 \oarg{副标题}
% 就排。写死 \texttt{true} 或 \texttt{false} 可以强制。
%    \begin{macrocode}
\tl_new:N \g__hit_subtitle_tl  \tl_gset:Nn \g__hit_subtitle_tl { auto }
\keys_define:nn { hit / class }
  { subtitle .code:n = { \__hit_tristate_set:Nn \g__hit_subtitle_tl {#1} } ,
    subtitle .default:n = { true } , }
%    \end{macrocode}
%
% \changes{v3.2a}{0000/00/00}{增加选项：哈工深官方修改，提供切换到毕业设计的选项}
%    \begin{macrocode}
\hit@declare@boolean@option{design}{false}
%    \end{macrocode}
%
% 是否开启debug模式（默认否）。如果开启，载入显示行号等的包，只为开发调试用。
%    \begin{macrocode}
\hit@declare@boolean@option{debug}{false}
%    \end{macrocode}
%
% \changes{v2.0.0}{2018/06/14}{此处删除newgeometry选项}
% 是否使用右开页（默认否）。
%    \begin{macrocode}
\hit@declare@boolean@option{openright}{false}
%    \end{macrocode}
%
% \changes{v2.0.10}{2019/06/25}{此处添加是否为提交图书馆电子版}
% 是否为提交图书馆电子版。
%    \begin{macrocode}
\hit@declare@boolean@option{library}{false}
%    \end{macrocode}
%
% 图题和标题最后一行是否居中对齐（默认是，非规范要求）。
% \changes{v1.0.06}{2017/10/25}{此处更改了选项的名称}
%    \begin{macrocode}
\hit@declare@boolean@option{capcenterlast}{false}
%    \end{macrocode}
%
% 子图图题和标题最后一行是否居中对齐（默认是，非规范要求）。
% \changes{v1.0.06}{2017/10/25}{此处添加子图最后一行图题是否居中选项}
%    \begin{macrocode}
\hit@declare@boolean@option{subcapcenterlast}{false}
%    \end{macrocode}
%
% 中文目录中Abstract是否均为大写
% \changes{v1.0.13}{2018/04/05}{此处添加中文目录中Abstract是否均为大写选项}
%    \begin{macrocode}
\hit@declare@boolean@option{absupper}{false}
%    \end{macrocode}
%
% 此处添加控制本科论文的页码横线选项
% \changes{v1.0.15}{2018/06/05}{添加控制本科论文的页码横线选项}
% 由于本科论文模板的页码格式现已统一，修改bsmainpagenumberline默认设置为true
% \changes{v3.0.16}{2021/11/29}{由于本科论文模板的页码格式现已统一，修改bsmainpagenumberline默认设置为true}
%    \begin{macrocode}
\hit@declare@boolean@option{bsmainpagenumberline}{true}
%    \end{macrocode}
%
% \changes{v3.1d}{2025/03/03}{哈尔滨本科要求 \cmd{\frontmatter} 中的页脚也要有横线，希望以后会统一。\issue[172]}
% \changes{v3.2a}{0000/00/00}{修正 \option{bsfrontpagenumberline} 的条件默认值。
%   原来把嵌套条件写在选项声明处，而那里在 \cs{ProcessKeyOptions} 之前执行，
%   此时 campus 与 type 尚未解析、\option{harbin} 与 \option{bachelor} 被兜底逻辑置真，
%   条件恒取第一支，任何组合下默认值都是 true。改为统一声明 false，选项解析完成
%   后再按校区与学位补默认值，并记录用户是否显式设置过以免覆盖。}
% \changes{v3.2a}{0000/00/00}{条件默认值改由三态取值承载，不再需要
%   记录用户是否显式设过。}
% 这个选项的默认值随校区与学位而变，交给三态取值处理：\texttt{auto}
% 时按哈尔滨本科为真解析，用户写了具体值就用具体值。
%    \begin{macrocode}
\hit@declare@boolean@option{bsfrontpagenumberline}{false}
\hit@declare@boolean@option{bsheadrule}{true}
%    \end{macrocode}
%
% 数学字体是否使用新罗马
% \changes{v2.0.05}{2018/12/05}{添加数学字体开关}
% \changes{v3.2a}{0000/00/00}{数学引擎换成 \pkg{unicode-math}，\option{newtxmath}
%   降级为 \option{math-font} 的一档，默认值改为假}
%    \begin{macrocode}
\hit@declare@boolean@option{newtxmath}{false}
%    \end{macrocode}
%
% \changes{v3.2a}{0000/00/00}{新增 \option{math-font} 与 \option{math-style} 选项}
% 数学字体与数学排版风格。
%    \begin{macrocode}
\hit@declare@string@option@as{math-font}{mathfont}
\hit@declare@string@option@as{math-style}{mathstyle}
%    \end{macrocode}
%
% 此处应广大刀客要求添加一参考文献分割开关
% \changes{v2.0.03}{2018/10/08}{添加参考文献分割开关}
%    \begin{macrocode}
\hit@declare@boolean@option{splitbibitem}{false}
%    \end{macrocode}
%
% 此处添加英文目录开关
% \changes{v3.0.11}{2020/10/29}{此处添加英文目录开关}
%    \begin{macrocode}
\hit@declare@boolean@option{engtoc}{false}
%    \end{macrocode}
%
% 声明字体选项。
%    \begin{macrocode}
\hit@declare@string@option{fontset}
%    \end{macrocode}
%
% \changes{v3.2a}{0000/00/00}{西文字体独立成 \option{font} 选项}
% 西文字体。原来一律用 TeX Gyre Termes（Times 的克隆），但学校要求的是
% Times New Roman，Windows 与 macOS 上本来就有。\option{auto} 按 \option{fontset}
% 决定：\texttt{windows}、\texttt{mac} 用 Times New Roman，其余用 TeX Gyre Termes。
%    \begin{macrocode}
\hit@declare@string@option@as{font-en}{font@en}
%    \end{macrocode}
%
% \changes{v3.2a}{0000/00/00}{中文字体独立成 \option{cjk-font} 选项}
% 中文字体。不写就照旧把 \option{fontset} 转给 \pkg{ctexbook}，由 \pkg{ctex} 设置；
% 写了就由模板自己设，\pkg{ctexbook} 收到的是 \texttt{fontset=none}。这样默认行为
% 一点没变，想换字体来源的人才走新路。
%    \begin{macrocode}
\hit@declare@string@option@as{font-zh}{font@zh}
%    \end{macrocode}
%
% \changes{v3.2a}{0000/00/00}{\option{font} 与 \option{cjk-font} 改名
%   \option{font-en} 与 \option{font-zh}，与模板其余选项的中英后缀一致。旧名仍
%   生效，只在日志里提示}
%    \begin{macrocode}
\keys_define:nn { hit / class }
  {
    font     .tl_gset:c = { hit@font@en } ,
    cjk-font .tl_gset:c = { hit@font@zh } ,
  }
%    \end{macrocode}
%
% 将其余选项默认传递给 \pkg{ctexbook}。
%    \begin{macrocode}
\keys_define:nn { hit / class }
  {
    unknown .code:n = { \hit@pass@or@warn { ctexbook } {#1} } ,
  }
%    \end{macrocode}
%
% 解析转发到模块的选项。
%    \begin{macrocode}
\tl_new:N \g__hit_toc_english_tl  \tl_gset:Nn \g__hit_toc_english_tl { auto }
\tl_new:N \g__hit_toc_blank_tl    \tl_gset:Nn \g__hit_toc_blank_tl   { auto }
\tl_new:N \g__hit_toc_four_tl     \tl_gset:Nn \g__hit_toc_four_tl    { auto }
\tl_new:N \g__hit_page_number_line_main_tl   \tl_gset:Nn \g__hit_page_number_line_main_tl  { auto }
\tl_new:N \g__hit_page_number_line_front_tl  \tl_gset:Nn \g__hit_page_number_line_front_tl { auto }
\tl_new:N \g__hit_head_rule_tl      \tl_gset:Nn \g__hit_head_rule_tl     { auto }
\keys_define:nn { hit / class }
  {
    tocfour  .code:n = { \tl_gset:Nn \g__hit_toc_four_tl    {#1} } ,
    tocblank .code:n = { \tl_gset:Nn \g__hit_toc_blank_tl   {#1} } ,
    engtoc   .code:n = { \tl_gset:Nn \g__hit_toc_english_tl {#1} } ,
    bsmainpagenumberline  .code:n = { \tl_gset:Nn \g__hit_page_number_line_main_tl  {#1} } ,
    bsfrontpagenumberline .code:n = { \tl_gset:Nn \g__hit_page_number_line_front_tl {#1} } ,
    bsheadrule            .code:n = { \tl_gset:Nn \g__hit_head_rule_tl     {#1} } ,
    tocfour  .default:n = { true } , tocblank .default:n = { true } ,
    engtoc   .default:n = { true } ,
    bsmainpagenumberline  .default:n = { true } ,
    bsfrontpagenumberline .default:n = { true } ,
    bsheadrule            .default:n = { true } ,
  }
\ProcessKeyOptions [~hit~/~class~]
%    \end{macrocode}
%
% \changes{v3.2a}{0000/00/00}{深圳校区本科现在直接指向本部本科模板，旧深圳本科模板代码暂时保留}
% \changes{v3.2a}{0000/00/00}{选项后处理（shenzhen重定向、type校验、PassOptions、fontset/zijv/geometrynew 默认）从 bookcls 移入 book-options，cls 只保留 LoadClass}
% 深圳校区本科现在直接指向本部本科模板，旧深圳本科模板代码暂时保留。
%    \begin{macrocode}
\bool_lazy_and:nnT
  { \bool_if_p:N \g__hit_shenzhen_bool }
  { \bool_if_p:N \g__hit_bachelor_bool }
  {
    \bool_gset_false:N \g__hit_shenzhen_bool
    \bool_gset_true:N \g__hit_harbin_bool
  }
%    \end{macrocode}
%
% 校区与学位这时才定下来，\option{bsfrontpagenumberline} 的默认值在此补齐。
% 用户显式设过就不动。
%    \begin{macrocode}

%    \end{macrocode}
%
% 用户至少要提供一个选项，指定论文类型。
% \changes{v3.0.0}{2020/05/20}{此处添加博后选项}
%
% 这条校验目前触发不到。选项解析之前有一段兜底，在 \option{master}、
% \option{doctor}、\option{postdoc} 都未置位时把 \option{bachelor} 置真，
% 而那段兜底在 \cs{ProcessKeyOptions} 之前执行，所以到这里
% \option{bachelor} 一定为真，四个条件不可能同时落空。实测
% \cs{documentclass}\oarg{campus=harbin} 不给 \option{type} 时不报错，
% 直接按本科排版。\file{hithesisart.cls} 没有这段兜底，它的同名校验能正常
% 触发。改动会让不写 \option{type} 的文档从「默认本科」变成报错，属用户
% 可见变更，此处只记录。
%    \begin{macrocode}
\cs_new_protected:Npn \hit@thesis@error@type
  {
    \ClassError { hithesis }
      {
        \MessageBreak Please~specify~thesis~type~in~option:~
        \MessageBreak type=[bachelor~|~master~|~doctor~|~postdoc]
      }
      { }
  }
\bool_lazy_any:nF
  {
    { \bool_if_p:N \g__hit_bachelor_bool }
    { \bool_if_p:N \g__hit_master_bool }
    { \bool_if_p:N \g__hit_doctor_bool }
    { \bool_if_p:N \g__hit_postdoc_bool }
  }
  { \hit@thesis@error@type }
%    \end{macrocode}
%
% 使用 \XeTeX\ 引擎时，\pkg{fontspec} 宏包会被 \pkg{xeCJK} 自动调用。传递
% 给 \pkg{fontspec} 宏包 \option{no-math} 选项，避免部分数学符号字体自动调整
% 为 CMR。其他引擎下没有这个问题，这一行会被无视。
%    \begin{macrocode}
\PassOptionsToPackage{no-math}{fontspec}
%    \end{macrocode}
%
% \pkg{newtx} 更新至 v1.7 版本后与 \pkg{ctex} 宏包冲突, 需要手动引入 \option{nofontspec} 选项. 判断方法为是否存在 v1.7 更新后的新宏包 \pkg{newtx.sty}
% \changes{v3.0.17}{2022/02/23}{\pkg{newtx} 更新至 v1.7 版本后与 \pkg{ctex} 宏包冲突, 需要手动引入 \option{nofontspec} 选项. 判断方法为是否存在 v1.7 更新后的新宏包 \pkg{newtx.sty}}
% 载入单双面打印设置，本、硕单面，博士双面。
%    \begin{macrocode}
\bool_lazy_or:nnT
  { \bool_if_p:N \g__hit_bachelor_bool }
  { \bool_if_p:N \g__hit_master_bool }
  { \PassOptionsToClass { oneside } { book } }
\bool_if:NT \g__hit_doctor_bool { \PassOptionsToClass { twoside } { book } }
%    \end{macrocode}
%
% \changes{v1.0.02}{2017/08/27}{添加了思源字体说明}
% \changes{v3.2a}{0000/00/00}{去掉模板自带的 siyuan fontset}
% 设置字体。宋体没有粗体，而标题要求粗宋，于是面临 \CTeX\ 那个经典的伪粗体 bug：
% 「首次出现伪粗体字体之后的正常字体无法复制」。自带粗宋的字体不必走伪粗体，
% 模板早年为此提供过一个 \texttt{siyuan} fontset，现在 \pkg{ctex} 的
% \texttt{ubuntu} 用的 Noto CJK 就是同一套思源，这一档已经去掉。
%
% \changes{v3.2a}{0000/00/00}{类内部不再用 \cs{ifthenelse}，\pkg{ifthen} 只为
%   兼容旧文档而保留}
% \pkg{ifthen} 模板自己已经不用了（三处 \cs{ifthenelse} 换成了 \cs{hit@if@same@TF}），
% 留着是怕旧论文里直接写了 \cs{ifthenelse} 而没自己加载。
%    \begin{macrocode}
\RequirePackage{ifthen}
\tl_if_empty:NTF \hit@font@zh
  {
    \PassOptionsToPackage { AutoFakeBold=2 } { xeCJK }
    \tl_if_empty:NF \hit@fontset
      { \PassOptionsToClass { fontset=\hit@fontset } { ctexbook } }
  }
  { \PassOptionsToClass { fontset=none } { ctexbook } }
%    \end{macrocode}
%
% 设置字距默认选项
% \changes{v3.0.13}{2020/11/10}{设置字距默认选项}
%    \begin{macrocode}
\bool_lazy_or:nnF
  { \bool_if_p:N \g__hit_ziju_regu_bool }
  { \bool_if_p:N \g__hit_ziju_word_bool }
  { \bool_gset_true:N \g__hit_ziju_regu_bool }
%    \end{macrocode}
%
% 重新设置geometry的默认选项
% \changes{v3.0.13}{2020/11/10}{重新设置geometry的默认选项}
%    \begin{macrocode}
\bool_lazy_any:nF
  {
    { \bool_if_p:N \g__hit_newgeometry_one_bool }
    { \bool_if_p:N \g__hit_newgeometry_two_bool }
    { \bool_if_p:N \g__hit_newgeometry_no_bool }
  }
  { \bool_gset_true:N \g__hit_newgeometry_two_bool }
%    \end{macrocode}
%
% \changes{v3.2a}{0000/00/00}{\option{style/glue} 决定的伸缩量抽成
%   \cs{hit@glue@skip} 与 \cs{hit@glue@font@size}，原地重复了二十多次的
%   \cs{ifhit@glue} 判断收敛到这两个函数里。}
% 读标志位的那族判断函数（\cs{hit@if@option@TF} 等）定义在
% \file{src/shared/hit-options.dtx}，各个类共用。按校区、学位组合出来的判断由
% \cs{hit@define@condition} 生成，一条一行：\cs{hit@if@harbin@master@or@doctor}、
% \cs{hit@if@harbin@master@or@doctor@or@shenzhen@master}、
% \cs{hit@if@shenzhen@master@or@harbin@bachelor}，以及决定要不要在封面上标出
% 学生类别的 \cs{hit@if@show@student@type}。名字摊开写，读条件本身就知道判的是什么。
% 参数从 expl3 之外传进来，里面的空格照原样保留。注意原写法里 \cs{fi} 之后的
% 空格会被记号化跳过，换成以 \texttt{\}} 结尾之后那个空格就成了真的空格，
% 所以替换时要把它一起吃掉。
%
% 标题前后的间距与行距。\option{style/glue} 开时给伸缩量，关时是固定值。
% 两个函数都可展开，所以能直接写在 \cs{ctexset} 的取值里。

%
% \changes{v3.2a}{0000/00/00}{加 \cs{hit@if@same@TF}，替掉三处 \cs{ifthenelse}}
% \cs{hit@if@same@TF}\marg{甲}\marg{乙}\marg{真}\marg{假} 比两段文字是否相同。不能用
% \cs{str\_if\_eq:eeTF}，那是完全展开后比较，遇上 \cs{hspace} 这类不可展开的命令
% 会留下悬空的 \cs{fi}；被比的正是章标题，\cs{cabstractcname} 在非本科档就是
% 「摘\cs{hspace}\marg{\cs{ccwd}}要」。这里照 \cs{ifthenelse}\marg{\cs{equal}…}
% 的做法先各自 \cs{protected@edef} 一遍，把 robust 命令挡住，再比记号表。
%    \begin{macrocode}
\tl_new:N \l__hit_compare_first_tl
\tl_new:N \l__hit_compare_second_tl
\hit@define@condition{harbin@master@or@doctor}{ \g__hit_harbin_bool && ( \g__hit_master_bool || \g__hit_doctor_bool ) }
\hit@define@condition{shenzhen@master@or@harbin@bachelor}{ ( \g__hit_shenzhen_bool && \g__hit_master_bool )
     || ( \g__hit_harbin_bool && \g__hit_bachelor_bool ) }
\hit@define@condition{harbin@master@or@doctor@or@shenzhen@master}{ ( \g__hit_harbin_bool && ( \g__hit_master_bool || \g__hit_doctor_bool ) )
     || ( \g__hit_shenzhen_bool && \g__hit_master_bool ) }
\hit@define@condition{show@student@type}{ ! ( ! \g__hit_shenzhen_bool && ! \g__hit_master_bool && \g__hit_fulltime_bool ) }
%    \end{macrocode}
%
%    \begin{macrocode}
%</thesis-options>
%<*thesiscls-load>
%    \end{macrocode}
%
% 使用 \pkg{ctexbook} 类，优于调用 \pkg{ctex} 宏包。
%    \begin{macrocode}
\LoadClass[a4paper,openany,UTF8,zihao=-4,scheme=plain]{ctexbook}
\bool_gset_true:N \g__hit_base_class_loaded_bool
%    \end{macrocode}
%</thesiscls-load>
%
% \changes{v3.2a}{0000/00/00}{模块内容改由 \file{hithesis.ins} 按运行期顺序直接拼入 cls，不再生成 \file{modules/*.sty}}
%
%<*thesiscls-tail>
%    \begin{macrocode}
\AtEndOfClass{\sloppy}
%    \end{macrocode}
%
% \changes{v3.2a}{0000/00/00}{重组 bookcls 结构：book-deps 仅保留第三方依赖与 hyperlink/geometry 基础设施，晚于它的功能模块（glossary/mainmatter/paragraph/footnote/math/floats/toc）提升至 cls；misc-book 模块合并为 cls 内联钩子}
%    \begin{macrocode}
%</thesiscls-tail>
%    \end{macrocode}
% ^^A ======================================================================
% ^^A Module thesis-deps: 第三方宏包集中加载与基础配置（含 hyperlink/geometry 内联）（hit-thesis）
% ^^A ======================================================================
%<*thesis-deps-a>
%
%
% \subsection{装载宏包}
% \label{sec:loadpackage}
%
% 引用的宏包和相应的定义。
%    \begin{macrocode}
\RequirePackage{etoolbox}
%    \end{macrocode}
%
% \changes{v3.2a}{0000/00/00}{book-deps-a 的引擎检查改用 \cs{sys\_if\_engine\_xetex:F}，
%   与 art-deps-a 保持一致。报错文本里的空格写成 \texttt{\~{}}。}
%    \begin{macrocode}
\sys_if_engine_xetex:F
  { \ClassError { hithesis } { Please~use:~\MessageBreak~xelatex } { } }
\RequirePackage{xparse}
%    \end{macrocode}
%
% \changes{v3.2a}{0000/00/00}{数学部分抽到 \file{src/shared/hit-math.dtx}}
% 数学宏包、定理机制与数学字体的设置见 \file{src/shared/hit-math.dtx}。
%    \begin{macrocode}
\hit@setup@math
%    \end{macrocode}
%
% 删除 \pkg{newtxtext}，由于最新一次更新 (\texttt{revision 62369}) 中与 \pkg{ctex} 宏包冲突，且无法解决，将正文字体替换为 TG TermesX 与 TG Heros.
% \changes{v3.0.18}{2022/05/02}{删除 \pkg{newtxtext}，由于最新一次更新 (\texttt{revision 62369}) 中与 \pkg{ctex} 宏包冲突，且无法解决，将正文字体替换为 TG Termes 与 TG Heros.}
% \changes{v3.0.19}{2022/05/05}{将 TG Termes 字体替换为 TG TermesX 字体，与 \pkg{newtxmath} 宏包更适配，删除 \option{Scale} 选项。}
% \changes{v3.2a}{0000/00/00}{英文字体三件套设置从 bookcls 拆出至 book-fonts 模块（与 art-fonts / artplus-fonts 保持一致）}
%    \begin{macrocode}
%</thesis-deps-a>
%    \end{macrocode}
% ^^A ======================================================================
% ^^A Module thesis-fonts: 正文字号、中英文主字体与字号宏命令（\wuhao 等）（hit-thesis）
% ^^A ======================================================================
%    \begin{macrocode}
%<*thesis-fonts>
%    \end{macrocode}
%
% \changes{v3.2a}{0000/00/00}{字体层主体抽到 \file{src/shared/hit-fonts.dtx}，
%   这里只留调用与本类特有的部分}
%    \begin{macrocode}
\hit@setup@fonts
\cs_set:Npn \normalsize
  {
    \@setfontsize \normalsize { 12bp } { \hit@glue@skip{20.50394bp}{2.834646bp} }
    \abovedisplayskip = 8pt
    \abovedisplayshortskip = 8pt
    \belowdisplayskip = \abovedisplayskip
    \belowdisplayshortskip = \abovedisplayshortskip
  }
%</thesis-fonts>
%<*thesis-deps-a>
%    \end{macrocode}
%
% 图形支持宏包。
%    \begin{macrocode}
\RequirePackage{graphicx}
%    \end{macrocode}
%
% \changes{v3.2a}{0000/00/00}{原 hithesis.sty 中的 xcolor、tikz 与全局颜色主题并入 book-deps，作为彩色与图形基础宏包}
% 彩色支持与 \pkg{tikz}（封面与正文图形需要），并预定义全局颜色主题。
%    \begin{macrocode}
\RequirePackage{xcolor}
\RequirePackage{tikz}
\definecolor{colorzero}{rgb}{0,~0,~0}
\definecolor{colorone}{rgb}{1,~0,~0}
\definecolor{colortwo}{rgb}{0,~0,~1}
\definecolor{colorthree}{rgb}{0,~1,~0}
%    \end{macrocode}
%
% \pkg{pdfpages} 宏包便于我们插入扫描后的授权页和声明页 PDF 文档。
%    \begin{macrocode}
\RequirePackage{pdfpages}
\includepdfset{fitpaper=true}
%    \end{macrocode}
%
% 更好的列表环境。
%    \begin{macrocode}
\hit@setup@list@packages
\RequirePackage{environ}
%    \end{macrocode}
%
% 禁止 \LaTeX 自动调整多余的页面底部空白，并保持脚注仍然在底部。
% 脚注按页编号。
%    \begin{macrocode}
\hit@setup@footnote@packages
%    \end{macrocode}
%
% 脚注格式。
%
% \changes{v3.2a}{0000/00/00}{删掉 \option{pifootnote} 与 \pkg{pifont}。带圈脚注号
%   统一由 \cs{quan} 排：先取当前宋体的真字符（\texttt{①}–\texttt{⑩}），超出字体
%   可用范围才画圈。\pkg{pifont} 那套 Dingbats 字形与正文中文字体不搭}
%
% 利用 \pkg{xeCJKfntef} 实现汉字的下划线和盒子内两段对齐，并可以避免
% \cs{makebox}\oarg{width}\oarg{s} 可能产生的 underful boxes。
%    \begin{macrocode}
\RequirePackage{xeCJKfntef}
%    \end{macrocode}
%
% 表格控制
%    \begin{macrocode}
\RequirePackage{longtable}
%    \end{macrocode}
%
% 使用三线表：\cs{toprule}，\cs{midrule}，\cs{bottomrule}。
%    \begin{macrocode}
\RequirePackage{booktabs}
%    \end{macrocode}
%
% 重定义 \pkg{booktabs} 宏包的默认线宽，绘制三线表时不再需要显式设置，以实现更好的内容与格式分离
% \changes{v3.0.20}{2022/05/06}{重定义 \pkg{booktabs} 宏包的默认线宽，绘制三线表时不再需要显式设置，以实现更好的内容与格式分离}
%    \begin{macrocode}
\heavyrulewidth=1.5pt
\lightrulewidth=1pt
\cmidrulewidth=0.5pt
%    \end{macrocode}
%
% 参考文献引用宏包。
% \changes{v3.0.15}{2020/05/21}{删去natbib package}
%    \begin{macrocode}
% \RequirePackage[sort&compress]{natbib}
%    \end{macrocode}
%
% 生成有书签的 pdf 及其开关，请结合 gbk2uni 避免书签乱码。
% \changes{v3.2a}{0000/00/00}{hyperref/url/gbt7714/citetwocomma 等链接与引用配置从 bookcls 拆出至 book-hyperlink 模块}
%    \begin{macrocode}
%</thesis-deps-a>
%    \end{macrocode}
% ^^A ======================================================================
% ^^A Module thesis-hyperlink: hyperref/url 配置（PDF 书签、链接颜色）（hit-thesis）
% ^^A ======================================================================
%    \begin{macrocode}
%<*thesis-hyperlink>
\hit@setup@hyperlink
\hit@setup@pdf@metadata
%    \end{macrocode}
%
% \changes{v3.2a}{0000/00/00}{删掉三处对 \pkg{natbib} 的 \cs{patchcmd}。
%   \cs{NAT@citexnum} 与 \cs{NAT@citex} 的参数文本是 \texttt{[\#1][\#2]\#3}，
%   带分隔符，\pkg{etoolbox} 的 \cs{patchcmd} 不支持这种宏，三处一直是失败的。
%   补丁想要的效果（页码注排在方括号外并上标）\pkg{gbt7714} 本来就提供了}
% \cs{NAT@citesuper} 是上标引用样式的排法，这个是真在起作用的，留着。
%
% \changes{v3.2a}{0000/00/00}{\option{citetwo} 的 comma 一档换个实现，从 2022 年加进来
%   就没生效过}
% \option{citetwo} 管相邻两个连续序号之间用什么连接：\texttt{endash} 排成
% \texttt{[1-2]}，\texttt{comma} 排成 \texttt{[1,2]}。三个及以上连续的照样压成
% \texttt{[1-3]}。
%
% 原先是拿 \cs{patchcmd} 去改 \cs{NAT@citexnum}，那个宏的参数文本是
% \texttt{[\#1][\#2]\#3}，带分隔符，\pkg{etoolbox} 补不了，所以一直是失败的。
% 改成动 \cs{def@NAT@last@yr}：它是 \pkg{natbib} 定义的普通宏，参数不带分隔符，
% 而且只在压缩连续序号的那一处被调用。区间刚开始时 \cs{NAT@last@yr} 还是
% \cs{relax}，这时改用普通分隔符 \cs{@citea}（就是逗号），后面几个仍旧走短横线。
% 判断放在运行期，所以写在 \cs{documentclass} 里还是导言区的 \cs{hitsetup} 里
% 都生效。
%    \begin{macrocode}
\cs_set_eq:NN \hit@original@natbib@last@year \def@NAT@last@yr
\cs_set:Npn \def@NAT@last@yr #1
  {
    \bool_lazy_and:nnTF
      { \bool_if_p:N \g__hit_citetwo_comma_bool }
      { \token_if_eq_meaning_p:NN \NAT@last@yr \relax }
      { \hit@original@natbib@last@year { \@citea } }
      { \hit@original@natbib@last@year {#1} }
  }
\cs_set:Npn \NAT@citesuper #1#2#3{\ifNAT@swa
  \if*#2*\else#2\NAT@spacechar\fi
\unskip\kern\p@\textsuperscript{\NAT@@open#1\NAT@@close\if*#3*\else#3\fi}%
   \else #1\fi\endgroup}
%    \end{macrocode}
%
% 根据窝工2021年的“四月革命”，引用文献标准更新到GBT2015
% \changes{v3.0.15}{2020/05/21}{添加参考文献GBT2015}
% \changes{v3.1c}{2023/04/16}{增加相邻两个参考文献连接符为逗号“,”的功能。\issue[137]}
%
% \changes{v3.2a}{0000/00/00}{\pkg{natbib} 的补丁提成一条命令，在 expl3 之外定义}
% \cs{patchcmd} 会把目标宏体 detokenize 后按当前 catcode 重扫，所以整条调用都得
% 在 expl3 之外写。而 \cs{ExplSyntaxOff} 塞进 \cs{bool\_if:NT} 的参数里是不生效的
% （参数早在读到 \cs{bool\_if:NT} 时就记号化完了），只能先在外面定义好，再由条件
% 调用。
%    \begin{macrocode}
%</thesis-hyperlink>
%<*thesis-deps-b>
%    \end{macrocode}
%
% \subsection{页面设置}
% \label{sec:layout}
% 本来这部分应该是最容易设置的，但根据我工\PGR\ 的3.8,3.4,3.2节的版芯矛盾，此处
% 设置两种版芯。
% \changes{v3.2a}{0000/00/00}{页面 geometry 设置从 bookcls 拆出至 book-geometry 模块}
%    \begin{macrocode}
%</thesis-deps-b>
%    \end{macrocode}
% ^^A ======================================================================
% ^^A Module thesis-geometry: 版芯尺寸与页边距（hit-thesis）
% ^^A ======================================================================
%    \begin{macrocode}
%<*thesis-geometry>
%    \end{macrocode}
%
% \changes{v3.2a}{0000/00/00}{book-geometry 模块整体改用 expl3 实现，是逐步迁移
%   的第一个试点。标志位当时仍由 kvoptions 声明（现已换成 expl3 布尔），
%   用 \cs{legacy\_if:nTF} 桥接。}
% 本模块是 expl3 迁移的试点。选项与标志位现在都由 \pkg{l3keys} 声明，
% 属于用户接口，本轮不动；这里只把模块内部的控制流换成 expl3 写法，
% 通过 \cs{legacy\_if:nTF} 读取旧标志位。
%    \begin{macrocode}
\bool_if:NTF \g__hit_debug_bool
  { \RequirePackage [ showframe ] { geometry } }
  { \RequirePackage { geometry } }
\geometry { a4paper , hcentering , ignoreall , nomarginpar }
%    \end{macrocode}
%
% 添加版芯设置选项
% \changes{v2.0.0}{2018/06/14}{添加版芯设置选项}
%    \begin{macrocode}
\bool_if:NTF \g__hit_newgeometry_two_bool
  {
    \geometry
      {
        centering ,
        text = { 150true~mm , 236true~mm } ,
        left = 30true~mm , head = 5true~mm , headsep = 2true~mm ,
        footskip = 0true~mm , foot = 5.2true~mm
      }
  }
  {
    \bool_if:NTF \g__hit_newgeometry_one_bool
      {
        \geometry
          {
            centering ,
            text = { 150true~mm , 240true~mm } ,
            left = 30true~mm , head = 5true~mm , headsep = 0true~mm ,
            footskip = 0true~mm , foot = 0true~mm
          }
      }
      {
        % 根据 PlutoThesis 原版定义而来
        \geometry
          {
            text = { 150true~mm , 224true~mm } ,
            top = 35.5true~mm , left = 30true~mm , head = 5true~mm ,
            headsep = 2.5true~mm , foot = 8.5true~mm
          }
      }
  }
%    \end{macrocode}
%
% \changes{v3.1d}{2025/03/03}{为哈尔滨本科特意改了页眉位置。\issue[172]}
% 为哈尔滨本科特意改了页眉位置。
%    \begin{macrocode}
\bool_lazy_and:nnT
  { \bool_if_p:N \g__hit_harbin_bool }
  { \bool_if_p:N \g__hit_bachelor_bool }
  {
    \geometry
      { top = 36.5true~mm , headsep = 1true~mm , bottom = 28.8true~mm }
  }
%</thesis-geometry>
%<*thesis-deps-c>
%    \end{macrocode}
%
% \changes{v3.2a}{0000/00/00}{book-deps-c 的两处条件加载改用 \cs{legacy\_if:nT}}
% 载入显示行号的包。
% \changes{v1.0.09}{2018/01/07}{添加debug包}
%    \begin{macrocode}
\bool_if:NT \g__hit_debug_bool
  {
    \RequirePackage { layout }
    \RequirePackage { layouts }
    \RequirePackage { lineno }
  }
%    \end{macrocode}
%
% 利用 \pkg{fancyhdr} 设置页眉页脚。
%    \begin{macrocode}
\RequirePackage{fancyhdr}
%    \end{macrocode}
%
% 其他包，表格、数学符号包
% \changes{v1.0.06}{2017/10/25}{此处添加子图最后一行图题是否居中选项}
%    \begin{macrocode}
\RequirePackage{tabularx}
\RequirePackage{xltabular}
\RequirePackage{varwidth}
%    \end{macrocode}
%
% 此处changepage环境用来控制索引页面的左右边距，规范中给出的示例的边距要大于正文。
% \changes{v1.0.10}{2018/02/19}{修改了索引的间距，使其更符合规范中的示例}
%    \begin{macrocode}
\RequirePackage{changepage}
\RequirePackage{multicol}
\RequirePackage[below]{placeins}%允许上一个section的浮动图形出现在下一个section的开始部分,还提供\FloatBarrier命令,使所有未处理的浮动图形立即被处理
\RequirePackage{flafter}       % 使得所有浮动体不能被放置在其浮动环境之前，以免浮动体在引述它的文本之前出现.
\RequirePackage{multirow}       %使用Multirow宏包，使得表格可以合并多个row格
\RequirePackage{subcaption}%支持子图，必须在 ccaption 之前
\RequirePackage{ccaption} %支持双语标题
%    \end{macrocode}
%
% \changes{v3.2a}{0000/00/00}{\cs{bicaption} 支持省略英文图表名的四参数写法}
% \pkg{ccaption} 的 \cs{bicaption} 是五个参数，第四个是第二条题注用的图表名，
% 用户得写成 \texttt{Fig.} 加一个负细空格才能把 \cs{fnum@figure} 插的空格抵掉。
% 这里让第四个参数可以省略，省略时取 \option{const} 族里的 \option{figure-prefix-en}
% 与 \option{table-prefix-en}。
%
% \changes{v3.2a}{0000/00/00}{五参数写法从报错改回可用，只提示一次。
%   v3.1e 的示例正文里就是五参数，报错会让老文档一行都编不过}
% 五个参数的写法照旧可用——那是 \pkg{ccaption} 的标准形式，v3.1e 的示例正文里
% 写的也是它，直接报错会让老文档编不过。现在只提示，v4.1 再移除。
%    \begin{macrocode}
\hit@setup@caption
\hit@setup@subfigure
%    \end{macrocode}
%
% 中英文索引包。
%    \begin{macrocode}
\RequirePackage[makeindex]{splitidx}
\newindex[]{china}
\newindex[]{english}
%    \end{macrocode}
%
% 排版logo。
% \changes{v3.0.21}{2022/05/11}{删除 \pkg{xltxtra} 宏包, 来实现规范样式的上标。\issue[125]}
%    \begin{macrocode}
%</thesis-deps-c>
%    \end{macrocode}
%
% \subsection{主文档格式}
% \label{sec:mainbody}
%
% \subsubsection{Three matters}
% \changes{v3.2a}{0000/00/00}{glossaries 与 frontmatter/mainmatter/backmatter 等设置从 bookcls 拆出至 book-glossary 与 book-mainmatter 模块}
% ^^A ======================================================================
% ^^A Module thesis-glossary: 术语与符号表（glossaries 配置 + 示例条目）（hit-thesis）
% ^^A ======================================================================
%    \begin{macrocode}
%<*thesis-glossary>
%    \end{macrocode}
%
% \changes{v3.1d}{2025/03/03}{将 \cs{glossarystyle} 替换为 \cs{setglossarystyle}。\issue[200]}
% \changes{v3.2a}{0000/00/00}{book-glossary 的两处 english 判断改用 \cs{legacy\_if:nTF}。
%   后面的 \cs{newacronym} 条目含英文字面文本，留在 expl3 之外。}
% 在 \pkg{glossaries} 中表示 \cs{glossarystle} 这个命令 Deprecated in v3.08a. Removed in v4.50.
%    \begin{macrocode}
\bool_if:NTF \g__hit_en_bool
  {
    \RequirePackage [ xindy , order=letter ] { glossaries-extra }
%    \end{macrocode}
%
% \changes{v3.2a}{0000/00/00}{去掉 \cs{newglossarystyle} 参数里那对不生效的
%   \cs{ExplSyntaxOn} 与 \cs{ExplSyntaxOff}}
% 样式代码存进 \pkg{glossaries} 时就已经记号化完了，写在里面的
% \cs{ExplSyntaxOn} 不起作用。环境正文本来就是排版内容，直接用 2e 写法。
%    \begin{macrocode}
\newglossarystyle{engstyle}{%
  \setglossarystyle{long}%
  \RenewDocumentEnvironment{theglossary}{}%
    {\vspace{-\intextsep}\centering
     \begin{longtable}{@{\hskip 1cm}p{3.5cm}p{\glsdescwidth}}}%
    {\end{longtable}}%
}
    \setabbreviationstyle [ acronym ] { long-short }
  }
  {
    \RequirePackage { glossaries }
    \setacronymstyle { short-long }
  }
\RenewDocumentCommand \genacrfullformat { m m } { \glsentrylong {#1} }
\bool_if:NTF \g__hit_en_bool { \makenoidxglossaries } { \makeglossaries }
%    \end{macrocode}
%
% \changes{v3.2a}{0000/00/00}{原 hithesis.sty 中的示例术语表条目与中文术语辅助宏（\cs{gtssbp}、\cs{gscna}、\cs{gscnas}）并入 book-glossary}
% 示例术语表条目（首次出现展开为全称，之后缩写），以及中文术语的双语索引辅助宏。
%    \begin{macrocode}
\newacronym{etssbp}{TSSBP}{Tree-structured~Stick-breaking~process}
\newacronym{etse}{TSE}{Taylor~Series~Expansion}
\newacronym{esvm}{SVM}{Support~Vector~Machine}
\newacronym{eml}{ML}{Machine~Learning}
\newacronym{eco}{CO}{Convex~Optimization}
\newacronym{tssbp}{树结构折筷过程}{树结构折筷过程（Tree-structured~Stick-breaking~process）}
\cs_new:Npn \gtssbp
  {
    \gls { tssbp } \sindex [china] { shu!树结构折筷过程 }
    \sindex [english] { Tree-structured~Stick-breaking~process }
  }
\newacronym[shortplural=SCNAs,longplural={体细胞拷贝数变异（Somatic~copy~number~alternation，SCNA）}]{scna}{SCNA}{体细胞拷贝数变异（Somatic~copy~number~alternation，SCNA）}
\cs_new:Npn \gscna
  {
    \gls { scna } \sindex [china] { ti!体细胞拷贝数变异 }
    \sindex [english] { Somatic~copy~number~alternation } \ignorespaces
  }
\cs_new:Npn \gscnas
  {
    \glspl { scna } \sindex [china] { ti!体细胞拷贝数变异 }
    \sindex [english] { Somatic~copy~number~alternation } \ignorespaces
  }
%</thesis-glossary>
%    \end{macrocode}
% ^^A ======================================================================
% ^^A Module thesis-mainmatter: 正文与前后文切换（frontmatter/mainmatter/backmatter）（hit-thesis）
% ^^A ======================================================================
%    \begin{macrocode}
%<*thesis-mainmatter>
%    \end{macrocode}
% \begin{macro}{\cleardoublepage}
% 对于 \textsl{openright} 选项，必须保证章首页右开，且如果前章末页无内容须
% 清空其页眉页脚。
% 如果\textsl{library}为真，则强制设置\textsl{openright}为真。
% \changes{v2.0.10}{2019/06/25}{添加\textsl{openright}和\textsl{library}逻辑}
% \changes{v3.0.0}{2020/05/20}{将glossaries移动到这里}
% \changes{v3.2a}{0000/00/00}{book-mainmatter 模块改用 expl3 实现。标志位仍由
%   kvoptions 声明，读用 \cs{legacy\_if:nTF}，写用 \cs{legacy\_if\_set\_false:n}（现已换成 expl3 布尔）。}
%    \begin{macrocode}
\bool_if:NT \g__hit_library_bool { \bool_gset_false:N \g__hit_openright_bool }
\cs_set_eq:NN \hit@clear@double@page \cleardoublepage
\cs_new:Npn \hit@clear@empty@double@page
  { \clearpage { \pagestyle { hit@empty } \hit@clear@double@page } }
\cs_set_eq:NN \cleardoublepage \hit@clear@empty@double@page
%    \end{macrocode}
%
% \end{macro}
% \begin{macro}{\frontmatter}
% 我们的单面和双面模式与常规的不太一样。
%    \begin{macrocode}
\cs_set:Npn \frontmatter
  {
    \bool_if:NTF \g__hit_openright_bool { \cleardoublepage } { \clearpage }
    \@mainmatterfalse
    \pagenumbering { Roman }
    \pagestyle { hit@empty }
%    \end{macrocode}
%
% 窝工深圳模板范例，一行34个字
% \changes{v3.0.11}{2020/10/29}{窝工深圳模板范例，一行34个字}
% 添加33行字选项
% \changes{v3.0.13}{2020/11/10}{添加33行字选项}
%    \begin{macrocode}
    \bool_if:NTF \g__hit_openright_bool { \cleardoublepage } { \clearpage }
    \bool_if:NTF \g__hit_ziju_word_bool
      { \ziju { 0.05416667 } }
      { \bool_if:NT \g__hit_ziju_regu_bool { \ziju { 0.05580808 } } }
  }
%    \end{macrocode}
%
% \end{macro}
% \begin{macro}{\mainmatter}
% 根据打印店（伪官方）的猛虎式操作，\cs{mainmatter}命令的逻辑是，双面打印时第一章必须在奇数页
% （不看文档别怪我）。
% \changes{v2.0.11}{2018/06/27}{设置第一章必须在奇数页}
%    \begin{macrocode}
\cs_set:Npn \mainmatter
  {
%    \end{macrocode}
% 规范中并没有加空行这一要求，此处不应该加，所以只在 \option{tocblank} 下才加。
%    \begin{macrocode}
    \__hit_tristate_if:NnT \g__hit_toc_blank_tl
      { ! \bool_if_p:N \g__hit_en_bool }
      {
        \addtocontents { toc } { \vspace { \baselineskip } }
        \addtocontents { toe } { \vspace { \baselineskip } }
      }
%    \end{macrocode}
%
% 添加英文版限制
% \changes{v3.0.0}{2020/05/21}{添加英文版限制}
%    \begin{macrocode}
    \bool_if:NTF \g__hit_doctor_bool
      { \bool_if:NTF \g__hit_library_bool { \clearpage } { \cleardoublepage } }
      { \clearpage }
    \@mainmattertrue
    \pagenumbering { arabic }
    \pagestyle { hit@headings }
  }
%    \end{macrocode}
%
% \end{macro}
% \begin{macro}{\backmatter}
%    \begin{macrocode}
\cs_set:Npn \backmatter
  {
    \bool_if:NTF \g__hit_openright_bool { \cleardoublepage } { \clearpage }
    \@mainmattertrue
  }
%</thesis-mainmatter>
%    \end{macrocode}
%
% \end{macro}
%
% \subsubsection{字体}
% \label{sec:font}
% \begin{macro}{\normalsize}
% 根据我工规定，正文小四号 (12bp) 字，行距为固定值3--4mm。
%    \begin{macrocode}

%    \end{macrocode}
%
% \end{macro}
% \subsubsection{页眉页脚}
% \label{sec:headerfooter}
% \begin{macro}{\hit@empty}
% \begin{macro}{\hit@plain}
% \begin{macro}{\hit@headings}
% 定义四种页眉页脚格式：
% \begin{itemize}
% \item \texttt{hit@empty}：页眉页脚都没有
% \item \texttt{hit@plain}：只显示页脚的页码。\cs{chapter} 自动调用
% \cs{thispagestyle\{hit@plain\}}。
% \item \texttt{hit@headings}：页眉页脚同时显示
% \item \texttt{hit@scanpage}：只显示页脚的页码，扫描页专用
% \end{itemize}
% ^^A ======================================================================
% ^^A Module thesis-pagestyle: 页眉页脚样式（fancyhdr 配置）（hit-thesis）
% ^^A ======================================================================
%    \begin{macrocode}
%<*thesis-pagestyle>
% \changes{v3.2a}{0000/00/00}{pagestyle 模块抽取后，fancyhdr 需提前加载}
% \changes{v3.2a}{0000/00/00}{book-pagestyle 的条件树改用 expl3。页眉横线与本科
%   页脚这两块内容留在 expl3 之外定义成宏：它们的花括号后换行会产生一个空格
%   记号，\cs{ExplSyntaxOn} 下会消失，见 CONTRIBUTING §4.6。}
\RequirePackage{fancyhdr}
\cs_set_eq:NN \hit@original@head@rule \headrule
\fancypagestyle{hit@empty}{%
  \fancyhf{}
  \cs_set_eq:NN \headrule \hit@original@head@rule
  \cs_set:Npn \headrulewidth {0pt}
  \cs_set:Npn \footrulewidth {0pt}
}
%    \end{macrocode}
%
% \changes{v3.2a}{0000/00/00}{schrb: 添加扫描PDF专用页面样式，只显示页码}
%    \begin{macrocode}
\fancypagestyle{hit@scanpage}{%
  \fancyhf{}
  \fancyfoot[C]{\xiaowu-\nobreakspace{}\thepage\nobreakspace{}-}
  \cs_set_eq:NN \headrule \relax
  \cs_set:Npn \headrulewidth {0pt}
  \cs_set:Npn \footrulewidth {0pt}
}
%    \end{macrocode}
%
% 页眉那条粗细双线三处都一样，本科页脚的页码格式只此一处。两块都含花括号后
% 换行产生的空格记号，所以留在 expl3 之外原样定义。
%    \begin{macrocode}
\cs_new_protected:Npn \hit@thesis@head@rule
  {
    \cs_set:Npn \headrule
      {
        \vskip 1.190132pt
        \hrule \@height 2.276208pt \@width \headwidth
        \vskip 0.75pt
        \hrule \@height .75pt \@width \headwidth
      }
  }
\cs_new_protected:Npn \hit@page@number@plain { \thepage }
\cs_new_protected:Npn \hit@page@number@decorated
  {
    \hit@if@option@TF{bachelor}%
      {-\hspace{\hit@page@number@separator}\thepage\hspace{\hit@page@number@separator}-}%
      {-\nobreakspace{}\thepage\nobreakspace{}-}%
  }
%    \end{macrocode}
%
% \changes{v3.2a}{0000/00/00}{页码装饰与页眉横线的开关改为三态，本科与非本科
%   走同一套逻辑。装饰形式仍按学位区分：本科用 \cs{hspace}，其余用不断行空格。}
% \option{pagenumberline} 的 \texttt{auto}：正文一律带装饰；前置部分除深圳与
% 威海本科外都带。\option{headrule} 的 \texttt{auto} 一律为真。
%    \begin{macrocode}
\cs_new_protected:Npn \hit@page@number@main
  {
    \__hit_tristate_if:NnTF \g__hit_page_number_line_main_tl { \c_true_bool }
      { \hit@page@number@decorated } { \hit@page@number@plain }
  }
\cs_new_protected:Npn \hit@page@number@front
  {
    \__hit_tristate_if:NnTF \g__hit_page_number_line_front_tl
      {
        ! \bool_lazy_and_p:nn
            { \bool_if_p:N \g__hit_bachelor_bool }
            {
              \bool_lazy_or_p:nn
                { \bool_if_p:N \g__hit_shenzhen_bool }
                { \bool_if_p:N \g__hit_weihai_bool }
            }
      }
      { \hit@page@number@decorated } { \hit@page@number@plain }
  }
\cs_new_protected:Npn \hit@thesis@head@rule@setup
  {
    \__hit_tristate_if:NnTF \g__hit_head_rule_tl { \c_true_bool }
      { \hit@thesis@head@rule } { \cs_set:Npn \headrulewidth { 0pt } }
  }
\cs_new_protected:Npn \hit@thesis@footer
  {
    \fancyfoot [C]
      {
        \xiaowu
        \legacy_if:nTF { @mainmatter }
          { \hit@page@number@main } { \hit@page@number@front }
      }
  }
%    \end{macrocode}
% 
% \changes{v3.1d}{2025/03/03}{定义页脚中横线与页码的距离。\issue[172]}
% 哈尔滨本科页脚的横线与页码的间距比原来的大，用判断来控制。
%    \begin{macrocode}
\skip_new:N \hit@page@number@separator
\bool_lazy_and:nnTF
  { \bool_if_p:N \g__hit_harbin_bool }
  { \bool_if_p:N \g__hit_bachelor_bool }
  { \dim_set:Nn \hit@page@number@separator { .3333333em } }
  { \dim_set:Nn \hit@page@number@separator { .1666667em } }
%    \end{macrocode}
%
% 此处根据本科生模板的多种版本，提供选项自定义页码、页眉样式。
% \changes{v1.0.15}{2018/06/05}{添加控制本科论文的页码横线选项}
% \changes{v1.0.15}{2018/06/05}{删除冗余的页面格式}
% \changes{v3.0.0}{2020/05/21}{添加英文版页眉格式控制}
% \changes{v3.0.0}{2020/05/21}{添加博后页眉控制}
% \changes{v3.1b}{2022/06/06}{本部硕博页眉中不再包含学科。}
% \changes{v3.1d}{2025/03/03}{根据 \issue[222] 修改深圳博士的页眉}
% \changes{v3.0.10}{2020/10/28}{根据窝工2020年深圳校区新规定，mainmatter之前都用正文内容作为页眉，之后都用哈尔滨工业大学作为页眉}
% \changes{v3.0.01}{2020/06/06}{此处设置威海本科页眉与本部一样}
%    \begin{macrocode}
\fancypagestyle { hit@headings }
  {
    \fancyhf { }
    \bool_if:NTF \g__hit_en_bool
      {
        \fancyhead [ C ] { \xiaowu [ 0 ] \leftmark }
        \fancyfoot [ C ] { \xiaowu - \nobreakspace \thepage \nobreakspace - }
        \hit@thesis@head@rule
      }
      {
        \bool_if:NT \g__hit_postdoc_bool
          {
            \fancyhead [ CO ] { \songti \xiaowu [ 0 ] \leftmark }
            \fancyhead [ CE ]
              { \songti \xiaowu [ 0 ] \hit@school@zh \hit@postdoc@report@zh }
          }
        \bool_if:NTF \g__hit_doctor_bool
          {
            \fancyhead [ CO ] { \songti \xiaowu [ 0 ] \leftmark }
            \bool_if:NT \g__hit_harbin_bool
              {
                \fancyhead [ CE ]
                  { \songti \xiaowu [ 0 ] \hit@school@zh \hit@degree@level@zh \hit@document@type@zh }
              }
            \bool_if:NT \g__hit_shenzhen_bool
              {
                \fancyhead [ CE ]
                  {
                    \songti \xiaowu [ 0 ] \hit@school@zh
                    \hit@if@option@TF{doctor}{}{\hit@shenzhen@campus}
                    \hit@degree@of@discipline@format@zh \hit@document@type@zh
                  }
              }
            \bool_if:NT \g__hit_weihai_bool
              {
                \fancyhead [ CE ]
                  {
                    \songti \xiaowu [ 0 ] \hit@school@zh \hit@weihai@campus
                    \hit@degree@of@discipline@format@zh \hit@document@type@zh
                  }
              }
          }
          {
            \bool_if:NT \g__hit_master_bool
              {
                \bool_if:NT \g__hit_harbin_bool
                  {
                    \fancyhead [ C ]
                      { \songti \xiaowu [ 0 ] \hit@school@zh \hit@degree@level@zh \hit@document@type@zh }
                  }
                \bool_if:NT \g__hit_shenzhen_bool
                  {
                    \legacy_if:nTF { @mainmatter }
                      {
                        \fancyhead [ C ]
                          {
                            \songti \xiaowu [ 0 ] \hit@school@zh \hit@degree@of@discipline@format@zh
                            \hit@document@type@zh
                          }
                      }
                      { \fancyhead [ C ] { \xiaowu [ 0 ] \leftmark } }
                  }
                \bool_if:NT \g__hit_weihai_bool
                  {
                    \fancyhead [ C ]
                      {
                        \songti \xiaowu [ 0 ] \hit@school@zh \hit@weihai@campus
                        \hit@degree@of@discipline@format@zh \hit@document@type@zh
                      }
                  }
              }
          }
        \bool_if:NTF \g__hit_bachelor_bool
          {
            \bool_if:NT \g__hit_harbin_bool
              {
                \fancyhead [ C ]
                  {
                    \songti \xiaowu [ 0 ] \hit@school@zh \hit@bachelor@degree@level@zh
                    \hit@bachelor@document@type@zh
                  }
              }
            \bool_if:NT \g__hit_shenzhen_bool
              {
                \fancyhead [ C ]
                  {
                    \songti \xiaowu [ 0 ] \hit@school@zh \hit@shenzhen@campus
                    \hit@bachelor@degree@level@zh \hit@bachelor@document@type@zh
                  }
              }
%    \end{macrocode}
% 威海本科页眉与本部一样。里层这个 \option{bachelor} 判断在外层已经成立，
% 恒取真支，保留是为了与改动前逐字对应。
%    \begin{macrocode}
            \bool_if:NT \g__hit_weihai_bool
              {
                \bool_if:NTF \g__hit_bachelor_bool
                  {
                    \fancyhead [ C ]
                      {
                        \songti \xiaowu [ 0 ] \hit@school@zh \hit@bachelor@degree@level@zh
                        \hit@bachelor@document@type@zh
                      }
                  }
                  {
                    \fancyhead [ C ]
                      {
                        \songti \xiaowu [ 0 ] \hit@school@zh \hit@weihai@campus
                        \hit@bachelor@degree@level@zh \hit@bachelor@document@type@zh
                      }
                  }
              }
            \hit@thesis@footer
            \hit@thesis@head@rule@setup
          }
          {
            \hit@thesis@footer
            \hit@thesis@head@rule@setup
          }
      }
    \cs_set:Npn \footrulewidth { 0pt }
  }
%    \end{macrocode}
%
% 此处解决页眉经典 bug。
%
% \changes{v3.2a}{0000/00/00}{\cs{ExplSyntaxOn} 从 \cs{AtBeginDocument} 的参数里
%   挪到外面。写在参数里不生效，\cs{@mkboth} 的两个参数会带上花括号内的空格}
%    \begin{macrocode}
\cs_new_protected:Npn \__hit_pagestyle_at_begin:
  {
    \pagestyle { hit@empty }
    \RenewDocumentCommand \chaptermark { m }
      { \@mkboth { \CTEXthechapter \enspace ##1 } { } }
  }
\AtBeginDocument { \__hit_pagestyle_at_begin: }
%</thesis-pagestyle>
%    \end{macrocode}
%
% \end{macro}
% \end{macro}
% \end{macro}
%
% \subsubsection{段落}
% \label{sec:paragraph}
%
% 全文首行缩进 2 字符，标点符号用全角
% \changes{v3.2a}{0000/00/00}{段落相关代码从 bookcls 拆出至 book-paragraph 模块}
% ^^A ======================================================================
% ^^A Module thesis-paragraph: 段落首行缩进与标点处理（hit-thesis）
% ^^A ======================================================================
%    \begin{macrocode}
%<*thesis-paragraph>
\ctexset{%
  punct=quanjiao,
  space=auto,
  autoindent=true}
%    \end{macrocode}
%
% 调整默认列表环境间的距离，以符合中文习惯。机制在
% \file{src/shared/hit-list.dtx}，这里只给取值。
%    \begin{macrocode}
\bool_gset_true:N \g__hit_list_nosep_bool
\clist_gset:Nn \g__hit_list_itemize_clist
  { leftmargin=0em,itemsep=0em,partopsep=0em,parsep=0em,topsep=0em,itemindent=3em }
\clist_gset:Nn \g__hit_list_enumerate_clist
  { leftmargin=0em,itemsep=0em,partopsep=0em,parsep=0em,topsep=0em,itemindent=3.5em }
\hit@setup@list@spacing
%    \end{macrocode}
%
% \changes{v3.2a}{0000/00/00}{列表间距的三行迁进 \file{src/shared/hit-list.dtx}，
%   与 \env{publist} 同处；这里只留调用与正文首行缩进}
% 正文首行缩进。
%    \begin{macrocode}
\dim_set:Nn \parindent {2em}
%    \end{macrocode}
%
% \changes{v3.2a}{0000/00/00}{\cs{cmd}/\cs{cs} 排版命令名工具从 book-math 迁入 book-paragraph（归位至正文排版细节）}
% 排版命令名的辅助：\cs{cs} 在文本中打印 \cs{name} 字面。
%    \begin{macrocode}
\cs_set:Npn \cmd #1 { \cs { \exp_after:wN \cmd@to@cs \token_to_str:N #1 } }
\cs_set:Npn \cmd@to@cs #1#2 { \char \number `#2 \relax }
\DeclareRobustCommand\cs[1]{\texttt{\char`\\#1}}
%</thesis-paragraph>
%    \end{macrocode}
%
% \subsubsection{脚注}
% \label{sec:footnote}
% 脚注符合中文习惯，数字带圈。
% \changes{v3.2a}{0000/00/00}{脚注相关代码从 bookcls 拆出至 book-footnote 模块}
% ^^A ======================================================================
% ^^A Module thesis-footnote: 脚注样式与编号（hit-thesis）
% ^^A ======================================================================
%    \begin{macrocode}
%<*thesis-footnote>
%    \end{macrocode}
%
% \changes{v3.2a}{0000/00/00}{脚注整层挪到 \file{src/shared/hit-footnote.dtx}，各个类共用}
% 脚注样式见 \file{src/shared/hit-footnote.dtx}。
%    \begin{macrocode}
\hit@setup@footnote
%</thesis-footnote>
%    \end{macrocode}
%
% \subsubsection{数学相关}
% \label{sec:equation}
% 允许太长的公式断行、分页等。
% \changes{v3.2a}{0000/00/00}{数学相关代码从 bookcls 拆出至 book-math 模块}
% ^^A ======================================================================
% ^^A Module thesis-math: 数学公式断行与定理样式（含 siunitx/bm/mathrsfs 等）（hit-thesis）
% ^^A ======================================================================
%    \begin{macrocode}
%<*thesis-math>
%    \end{macrocode}
%
% \changes{v3.2a}{0000/00/00}{原 hithesis.sty 中的数学相关宏包与符号工具并入 book-math}
% 根据窝工规范对数字书写的规定（4 位及以上每 3 位空 1/4 汉字），加载 \pkg{siunitx} 并配置分组分隔。
% \pkg{bm} 处理数学公式中的黑斜体，\pkg{mathrsfs} 提供不同于 \cs{mathcal}/\cs{mathfrak} 的英文花体。
%    \begin{macrocode}
\RequirePackage{siunitx}
\sisetup{group-minimum-digits=4,~group-separator=~\hspace{0.25em}}
\sisetup{detect-weight,detect-mode,detect-family}
%    \end{macrocode}
%
% \changes{v3.2a}{0000/00/00}{走 \pkg{unicode-math} 时 \cs{bm} 改用 \cs{symbf}}
% \pkg{bm} 只在 8~位数学字体下加载。\pkg{unicode-math} 的符号是扩展 mathchar，
% \pkg{bm} 处理不了，会报 \texttt{Extended mathchar used as mathchar}，
% 这时把 \cs{bm} 接到 \pkg{unicode-math} 自己的 \cs{symbf} 上。
%    \begin{macrocode}
\bool_if:NTF \g__hit_unicodemath_bool
  { \cs_set_eq:NN \bm \symbf }
  { \RequirePackage { bm } }
\RequirePackage{mathrsfs}
%    \end{macrocode}
%
% 常用数学记号约定：向量、矩阵、集合、有向/无向边、网络与节点。
%    \begin{macrocode}
\NewDocumentCommand \theVector         { m }   { \bm {#1} }
\NewDocumentCommand \theMatrix         { m }   { \mathbb {#1} }
\NewDocumentCommand \theSet            { m }   { \mathcal {#1} }
\NewDocumentCommand \theDirected       { m }   { { \overrightarrow {#1} } }
\NewDocumentCommand \theUndirected     { m }   { { \overline {#1} } }
\NewDocumentCommand \theNetwork        { m }   { \mathscr {#1} }
\NewDocumentCommand \theNode           { m }   { { \text {#1} } }
\NewDocumentCommand \theDirectedEdge   { m m } { { \overrightarrow { {#1} {#2} } } }
\NewDocumentCommand \theUndirectedEdge { m m } { { \overline { {#1} {#2} } } }
%    \end{macrocode}
%
% \changes{v3.2a}{0000/00/00}{\cs{cmd}/\cs{cs} 排版命令名工具从 book-math 迁出至 book-paragraph}
% 以下为原模板的核心公式断行设置。
%    \begin{macrocode}
\allowdisplaybreaks[4]
\predisplaypenalty=0%公式之前可以换页，公式出现在页面顶部
\postdisplaypenalty=0
\cs_set:Npn \theequation
  { \hit@if@greater@T{\c@chapter}{0}{\thechapter-} \@arabic \c@equation }
%    \end{macrocode}
%
% \changes{v2.0.03}{2018/10/08}{设置公式前后随意断页}
% 公式距前后文的距离由 4 个参数控制，参见 \cs{normalsize} 的定义。
% 同时为了让 \pkg{amsmath} 的 \cs{tag*} 命令得到正确的格式，我们必须修改这些代
% 码。\cs{make@df@tag} 是定义 \cs{tag*} 和 \cs{tag} 内部命令的。
% \cs{make@df@tag@@} 处理 \cs{tag*}，我们就改它！
% \begin{latex}
% \def\make@df@tag{\@ifstar\make@df@tag@@\make@df@tag@@@}
% \def\make@df@tag@@#1{%
%   \gdef\df@tag{\maketag@@@{#1}\def\@currentlabel{#1}}}
% \end{latex}
%    \begin{macrocode}
\cs_set:Npn \make@df@tag { \@ifstar \hit@make@starred@tag \make@df@tag@@@ }
\cs_set:Npn \hit@make@starred@tag #1
  { \cs_gset:Npn \df@tag { \hit@make@tag {#1} \cs_set:Npn \@currentlabel {#1} } }
\cs_set:Npn \hit@make@tag #1
  { \maketag@@@ { ( \ignorespaces #1 \unskip \@@italiccorr ) } }
\cs_set:Npn \tagform@ #1
  { \maketag@@@ { ( \ignorespaces #1 \unskip \@@italiccorr ) \equcaption {#1} } }
%    \end{macrocode}
%
% 修改 \cs{tagform} 会影响 \cs{eqref}。
%    \begin{macrocode}
\RenewDocumentCommand \eqref { m } { \textup { ( \ref {#1} ) } }
%    \end{macrocode}
%
% \changes{v3.2a}{0000/00/00}{ntheorem 全局样式（\cs{renewtheoremstyle}/\cs{theorembodyfont}/\cs{theoremheaderfont}/\cs{theoremsymbol}/\cs{theoremXXskipamount}）与 \cs{arraycolsep} 从 book-floats 迁入 book-math（归位至数学/定理职责）}
% ntheorem 定理全局样式与公式数组列间距。
%    \begin{macrocode}
\hit@theorem@style{
  body-font~~~=~{\songti\rmfamily},
  head-font~~~=~{\heiti\rmfamily},
  qed~~~~~~~~~=~{$\square$},
  note-braces~=~{}{},
  note-font~~~=~{},
  head-space~~=~{6pt},
  space-above~=~{1pt},
  space-below~=~{-2pt},
}
\arraycolsep=1.6pt
%    \end{macrocode}
%
% \changes{v3.2a}{0000/00/00}{\env{eqdenote} 的定义挪到 \file{src/shared/hit-math.dtx}}
% \env{eqdenote} 见 \file{src/shared/hit-math.dtx}。
%    \begin{macrocode}
%</thesis-math>
%    \end{macrocode}
%
% \subsubsection{浮动对象以及表格}
% \label{sec:float}
% 设置浮动对象和文字之间的距离，由于规范中没有明确规定，根据经验，设置成正文汉字
% 高度。
% \changes{v1.0.09}{2018/01/07}{修正float垂直间距bug}
% \changes{v2.0.09}{2019/06/24}{修正float垂直间距bug}
% \changes{v3.2a}{0000/00/00}{浮动对象与表格代码从 bookcls 拆出至 book-floats 模块；拼接位置留在 cls 中的原处以避免依赖错位}
% ^^A ======================================================================
% ^^A Module thesis-floats: 浮动体默认行为与双语题注（hit-thesis）
% ^^A ======================================================================
%    \begin{macrocode}
%<*thesis-floats>
%    \end{macrocode}
%
% \changes{v3.0.0}{2020/05/21}{添加博后的浮动间距设置}
% \changes{v3.2a}{0000/00/00}{浮动体间距的三个分支迁到 expl3。}
% 博后单独一套间距；其余按 \option{style/glue} 决定要不要给伸缩量。
%    \begin{macrocode}
\bool_if:NTF \g__hit_postdoc_bool
  {
    \skip_set:Nn \intextsep    { 10.5bp }
    \skip_set:Nn \textfloatsep { 10.5bp }
    \skip_set:Nn \floatsep     { 10.5bp }
  }
  {
    \bool_if:NTF \g__hit_glue_bool
      {
        \skip_set:Nn \intextsep    { 8.50398bp~plus~2.83465bp~minus~0bp }
        \skip_set:Nn \textfloatsep { 8.50398bp~plus~2.83465bp~minus~0bp }
        \skip_set:Nn \floatsep     { 12bp~plus~2.83465bp~minus~0bp }
      }
      {
        \skip_set:Nn \intextsep    { 8.50398bp }
        \skip_set:Nn \textfloatsep { 8.50398bp }
        \skip_set:Nn \floatsep     { 12bp }
      }
  }
%    \end{macrocode}
%
% 此处设置float在p选项时间隔，此处不设置\cs{@fptop}和\cs{@fpbot}以确保居中。
% \changes{v1.0.12}{2018/04/03}{修正float为p状态时默认不居中bug}
% \changes{v2.0.04}{2018/12/04}{删除\cs{@fpsep}设置，似乎没有什么用}
% \changes{v2.0.04}{2018/12/04}{更新\cs{intextsep}\cs{textfloatsep}\cs{floatsep}间距为正文行间距}
% 下面这组命令使浮动对象的缺省值稍微宽松一点，从而防止幅度对象占据过多的文本页面，
% 也可以防止在很大空白的浮动页上放置很小的图形。
% \changes{v1.0.08}{2017/11/05}{修改附录中图、表、公式数字编码}
%    \begin{macrocode}
\tl_gput_right:Nn \appendix { \cs_set_nopar:Npn \thefigure { \thechapter - \arabic {figure} } }
\tl_gput_right:Nn \appendix { \cs_set_nopar:Npn \thetable { \thechapter - \arabic {table} } }
\tl_gput_right:Nn \appendix { \cs_set_nopar:Npn \theequation { \thechapter - \arabic {equation} } }
\cs_set:Npn \textfraction      { 0.15 }
\cs_set:Npn \topfraction       { 0.85 }
\cs_set:Npn \bottomfraction    { 0.65 }
\cs_set:Npn \floatpagefraction { 0.60 }
%    \end{macrocode}
%
% 由于我工的双标题，导致标题之下多出一空白字符的距离，去除。
% \changes{v2.0.04}{2018/12/04}{更新图段后空白距离}
% \changes{v2.0.04}{2018/12/04}{删除表段后空白距离}
% \changes{v2.0.05}{2018/12/05}{删除图段后空白距离}
%    \begin{macro}{\@makecaption}
% 根据我工规范，本科和硕博的图题序号之后的空格不一样。
%    \begin{hitrgu}[\PGR][2.13.1]
% 每个图均应有图题（由图序和图名组成），图题不宜有标点符号，图名在图序之后空1个
% 半角字符排写。
%    \end{hitrgu}
%    \begin{hitrgu}[\UGR][2.13.1]
% 每个图均应有图题（由图序和图名组成），图题不宜有标点符号，图名在图序之后空1个
% 字符排写。
%    \end{hitrgu}
% 我工规范中没有明确规定是否标题是否居中对齐，这里给出一个居中选项自行调整。
% 注意，我工只规定：“居中书写”。此处不额外添加悬挂处理。
% \changes{v1.0.06}{2017/10/25}{此处更改了选项的名称}
% \changes{v1.0.07}{2017/11/04}{优化了最后一行居中算法，使其两边对齐、单词内部断行}
%
% \changes{v3.2a}{0000/00/00}{\cs{@makecaption} 与 \cs{LT@makecaption} 迁到 expl3}
% \cs{hit@caption@separator} 是题号与题名之间的间隔，本科一个汉字宽，其余是 \cs{enskip}。
% \cs{hit@caption@last@line} 把最后一行单独取出来居中，\option{capcenterlast} 开时用。
% 里面的盒子编号与分组照抄原来的写法，只是换成 expl3 的拼法。量宽仍用 \cs{sbox}
% 而不是 \cs{hbox\_set:Nn}：前者会套一层 \cs{color@begingroup}，盒子里的颜色
% 改动不会漏出去，后者没有。
%    \begin{macrocode}
\cs_new:Npn \hit@caption@separator { \hit@if@option@TF { bachelor } { \hskip \ccwd } { \enskip } }
\cs_new_protected:Npn \hit@caption@last@line #1
  {
    \vskip 6.3bp
    {
      \tex_setbox:D 0 = \tex_vbox:D {#1}
      \tex_setbox:D 1 = \tex_vbox:D
        {
          \tex_unvbox:D 0
          \tex_setbox:D 2 = \tex_lastbox:D
          \hbox_to_wd:nn { \textwidth }
            { \hfill \tex_unhcopy:D 2 \unskip \unskip \hfill }
        }
      \tex_unvbox:D 1
    }
  }
\cs_set:Npn \@makecaption #1#2
  {
    \vskip \abovecaptionskip
    \wuhao
    \sbox \@tempboxa { #1 \hit@caption@separator #2 }
    \dim_compare:nNnTF { \box_wd:N \@tempboxa } > { \hsize }
      {
        \hit@if@option@TF { capcenterlast }
          { \hit@caption@last@line { #1 \hit@caption@separator #2 } }
          { #1 \hit@caption@separator #2 }
        \par
      }
      {
        \tex_global:D \@minipagefalse
        \hb@xt@ \hsize { \hfil \box \@tempboxa \hfil }
      }
    \vskip \belowcaptionskip
  }
%    \end{macrocode}
%
%    \end{macro}
%    \begin{macro}{\LT@makecaption}
% 有关longtable的caption支持
%    \begin{macrocode}
\cs_set:Npn \LT@makecaption #1#2#3
  {
    \LT@mcol \LT@cols c
      {
        \hbox_to_zero:n
          {
            \hss
            \parbox [t] \LTcapwidth
              {
                \wuhao
                \sbox \@tempboxa { #1 { #2 \hit@caption@separator } #3 }
                \dim_compare:nNnTF { \box_wd:N \@tempboxa } > { \hsize }
                  {
                    \hit@if@option@TF { capcenterlast }
                      { \hit@caption@last@line { #1 { #2 \hit@caption@separator } #3 } }
                      { #1 { #2 \hit@caption@separator } #3 }
                    \par
                  }
                  {
                    \tex_global:D \@minipagefalse
                    \hbox_to_wd:nn { \hsize } { \hfil \box \@tempboxa \hfil }
                  }
                \endgraf \vskip \belowcaptionskip
              }
            \hss
          }
      }
  }
%    \end{macrocode}
%
%    \end{macro}
%    \begin{macro}{\longbionenumcaption}
% 长表格的双语标题是一个坑. 因为第一不能用浮动格式，只能用longtable包中的tabular
% ，这样表题只能使用表格中前两行来写。这样出现了一个问题是，中英表题的间距，标题
% 和表第一行间距，表格内部间距等多个变量的协调问题。这个问题只要使用tabular的形
% 式，就是无解的。唯一的方法就是把这些参数都给用户列出来。以下，第2，5参数为中英
% 双语标题内容，1，4为标题参数。6为中英标题间距，7为表题和表格间距。
%    \begin{macrocode}
\cs_set_nopar:Npn \longbionenumcaption #1 #2 #3 #4 #5 #6 #7
  {
%
\@if@contemptyarg{#1}{\caption{#2}}{\caption[#1]{#2}}%
\cs_gset_eq:NN \@cont@oldtablename \tablename
\cs_gset_nopar:Npn \tablename {#3}~
\cs_gset_eq:NN \LT@c@ption \@cont@LT@nonumintoc
\\[#6]~
\@if@contemptyarg{#4}{\caption{#5}}{\caption[#4]{#5}}%
\cs_gset_eq:NN \tablename \@cont@oldtablename
\cs_gset_eq:NN \LT@c@ption \@cont@oldLT@c@ption
\vspace{#7}
  }
%    \end{macrocode}
%
%    \end{macro}
%    \begin{macro}{\ltfontsize}
% \changes{v3.2a}{0000/00/00}{长表格字号机制抽到 \file{src/shared/hit-table.dtx}}
% 我们采用 \pkg{longtable} 来处理跨页的表格。默认字号压到五号，机制见
% \file{src/shared/hit-table.dtx}。
%    \begin{macrocode}
\hit@setup@table
%    \end{macrocode}
%
%    \end{macro}
% 图表名称及格式。
% \changes{v1.0.08}{2017/11/05}{删除冗余公式符号定义}
%    \begin{macrocode}
\cs_set:Npn \thesubtable { ( \alph { subtable } ) }
\cs_set:Npn \thefigure   { \arabic { chapter } - \arabic { figure } }% 使图编号为 7-1 的格式
%    \end{macrocode}
%
% \changes{v3.0.15}{2020/05/21}{根据窝工2021“四月革命”之精神，修改子图为双括号}
%    \begin{macrocode}
\cs_set:Npn \thesubfigure { ( \alph { subfigure } ) }% 子图编号为 a) 的格式
\cs_set:Npn \p@subfigure  { \thefigure \nobreakspace { } }% 子图引用为 7-1 a)
\cs_set:Npn \thetable     { \arabic { chapter } - \arabic { table } }% 表编号为 7-1
%    \end{macrocode}
%
% \changes{v3.2a}{0000/00/00}{\cs{setitemize}/\cs{setenumerate} 从 book-floats 迁出至 book-paragraph}
% 此处删除hang caption的设置
% \changes{v1.0.06}{2017/10/25}{删除caption hang 的默认设置，因为不在规范要求中}
%    \begin{macrocode}
\captionnamefont{\wuhao}
\captiontitlefont{\wuhao}
\cs_set:Npn \subcapsize { \wuhao }
%    \end{macrocode}
%
% \option{modern} 下子图题的字体、标签分隔与最后一行对齐交给 \pkg{caption}，
% 键名不同但取的是同一批值。
%    \begin{macrocode}
\DeclareCaptionFont { hit@subcap } { \subcapsize }
\DeclareCaptionLabelSeparator { hit@sub } { \hskip \subfiglabelskip }
\captionsetup [ sub ]
  { font = hit@subcap , labelformat = simple , labelsep = hit@sub , skip = 0pt }
\bool_if:NT \g__hit_subcapcenterlast_bool
  { \captionsetup [ sub ] { justification = centerlast } }
\dim_set:Nn \abovecaptionskip {0pt}%为了双标题之间的间距，不能设置
\dim_set:Nn \belowcaptionskip {0pt}
%    \end{macrocode}
%
% \changes{v3.2a}{0000/00/00}{\env{publist} 挪到 \file{src/shared/hit-publist.dtx}，各个类共用}
% \env{publist} 见 \file{src/shared/hit-publist.dtx}。
%    \begin{macrocode}
\hit@setup@publist
%    \end{macrocode}
%
% \changes{v2.0.01}{2018/06/28}{去除定理注释括号}
% \changes{v3.2a}{0000/00/00}{ntheorem 定理样式与 \cs{arraycolsep} 从 book-floats 迁出至 book-math；\cs{setlength}\{\cs{parindent}\}\{2em\} 迁出至 book-paragraph}
%
% \changes{v3.2a}{0000/00/00}{原 hithesis.sty 中的旋转/算法/代码段宏包并入 book-floats}
% 旋转图表所需的 \pkg{rotating}。
%    \begin{macrocode}
\RequirePackage{rotating}
%    \end{macrocode}
%
% 算法宏包。注意加载顺序为 float、hyperref、algorithm2e，否则无法生成算法列表。
% 算法标题居中、中英双语标题、按章节编号等设置参见原文档。
%    \begin{macrocode}
\RequirePackage[algoruled,linesnumbered,algochapter]{algorithm2e}
\SetAlCapSty{}
\NewDocumentCommand \foocaption { m }
  {~\cs_set:Npn \@algocf@pre@plainruled
    {
      \hrule height 1.5pt depth 0pt \kern \interspacetitleruled #1~\kern \interspacealgoruled
      \hrule height 1pt depth 0pt \kern \interspacetitleruled
    }~}
\cs_set:Npn \@algocf@post@ruled
  { \kern \interspacealgoruled \hrule height 1.5pt \relax }
\cs_new:Npn \algoenname { Algo. }
\newfloatlist[chapter]{algoen}{aen}{\listalgoenname}{\algoenname}
\newfixedcaption{\algoencaption}{algoen}
\cs_set:Npn \thealgoen { \thechapter - \arabic { algocf } }
\cs_set:Npn \@cftmakeaentitle
  { \chapter * { \listalgoenname \@mkboth { \listalgoenname } { \listalgoenname } } }
\cs_set:Npn \algorithmcfname { 算法 }
\dim_set:Nn \AlCapSkip {1.2ex}
\SetAlgoSkip{1pt}
\RenewDocumentCommand \algocf@captiontext { m m }
  { \wuhao #1 \algocf@typo \nobreakspace { } \AlCapFnt { } #2 }
\cs_if_exist:NTF \algocf@within
  {
    \cs_set:Npn \thealgocf
      { \use:c { the \algocf@within } - \@arabic \c@algocf }
  }
  { \cs_set:Npn \thealgocf { \@arabic \c@algocf } }
\RenewDocumentCommand \algocf@makecaption { m m }
  {
    \dim_add:Nn \hsize { \algomargin }
    \sbox \@tempboxa { \algocf@captiontext {#1} {#2} }
    \hskip .5\algomargin
    \parbox [t] { \hsize } { \centering \algocf@captiontext {#1} {#2} }
    \dim_add:Nn \hsize { -\algomargin }
  }
\NewDocumentCommand \AlgoBiCaption { m m }
  {
    \caption [#1]
      { \protect \dim_set:Nn \baselineskip { 1.5em } #1~\protect~\\ Algo.\thealgocf \nobreakspace { } #2 }
    \addcontentsline { aen } { algoen } { \protect \numberline { \thealgoen } {#2} }
  }
%    \end{macrocode}
%
% 代码段排版 \pkg{listings}（默认柔性列宽与自动断行）。
%    \begin{macrocode}
\RequirePackage{listings}
\lstset{
columns=flexible,
breaklines=true
}
%    \end{macrocode}
%
% \changes{v3.2a}{0000/00/00}{原 hithesis.sty 中的示例 tikz 样式与染色体示意图绘制宏并入 book-floats}
% 示例文档使用的 tikz 样式与染色体段落、细胞绘制宏。
%    \begin{macrocode}
\tikzset{maternal/.style={colorone}}
\tikzset{paternal/.style={colortwo}}
\tikzset{variant/.style={colorthree!80!colorzero}}
\tikzset{reference/.style={colorzero}}
\tikzset{aallele/.style={colorzero,rotate=90}}
\tikzset{ballele/.style={colorthree!80!colorzero,rotate=90}}
\tikzset{refseg/.style={colorzero,draw=colorzero,~opacity=0.2}}
\tikzset{mseg/.style={colorone,draw=colorone,~opacity=0.2}}
\tikzset{pseg/.style={colortwo,draw=colortwo,~opacity=0.2}}
\tikzset{vseg/.style={colorthree!80!colorzero,draw=colorthree!80!colorzero,~opacity=0.6}}
\tikzset{bncell/.style={draw=colorzero,opacity=0.2,line~width=2pt,~rounded~corners=1pt}}
\tikzset{btcell/.style={draw=colorone,opacity=0.6,~line~width=2pt,~rounded~corners=1pt}}
\tikzset{tncell/.style={colorzero,opacity=0.9}}
\tikzset{ttcell/.style={colorone,opacity=0.6}}
\tikzset{tscell/.style={colorzero}}
\tikzset{refcell/.style={colorzero}}
\tikzset{evolve/.style={->,draw=colortwo,opacity=0.3,line~width=1.5pt}}
\tikzset{fakeevolve/.style={->,draw=colorzero,opacity=0.3,line~width=1.5pt}}
\tikzset{refline/.style={dashed,draw=colorzero,line~width=1pt}}
\tikzset{tnline/.style={dashed,draw=colorzero,opacity=0.3,line~width=1pt}}
\NewDocumentCommand \gseg { m m m m m m m m m }
  {
%
~\pgfmathsetmacro{\sstartx}{#1}~
 \pgfmathsetmacro{\slengx}{#2}~
 \pgfmathsetmacro{\sy}{#3}~
 \pgfmathsetmacro{\sdy}{#4}~
 \pgfmathsetmacro{\sdx}{#5}~
 \pgfmathsetmacro{\sdxh}{#7}~
 \pgfmathsetmacro{\sdxt}{#8}~
 \fill[#6](\sstartx,\sy)--(\sstartx-\sdx,\sy+\sdy)--~
 (\slengx+\sstartx+1.5-\sdx,\sy+\sdy)--(\slengx+\sstartx+1.5,\sy)--~
 (\slengx+\sstartx+1.5-\sdx,\sy-\sdy)--(\sstartx-\sdx,\sy-\sdy)--cycle;~
 \draw[#9]~(\sstartx-\sdxh,\sy)~--~(\sstartx,~\sy);~
 \draw[#9]~(\slengx+\sstartx+1.5,~\sy)~--~(\slengx+\sstartx+1.5+\sdxt,\sy);~
  }
\NewDocumentCommand \gsegr { m m m m m m m m m }
  {
%
~\pgfmathsetmacro{\sstartx}{#1}~
 \pgfmathsetmacro{\slengx}{#2}~
 \pgfmathsetmacro{\sy}{#3}~
 \pgfmathsetmacro{\sdy}{#4}~
 \pgfmathsetmacro{\sdx}{#5}~
 \pgfmathsetmacro{\sdxh}{#7}~
 \pgfmathsetmacro{\sdxt}{#8}~
 \fill[#6](\sstartx-0.5,\sy)--(\sstartx+\sdx-0.5,\sy+\sdy)--~
 (\slengx+\sstartx+1.5+\sdx-0.5,\sy+\sdy)--(\slengx+\sstartx+1.5-0.5,\sy)--~
 (\slengx+\sstartx+1.5+\sdx-0.5,\sy-\sdy)--(\sstartx+\sdx-0.5,\sy-\sdy)--cycle;~
 \draw[#9]~(\sstartx-\sdxh-0.5,\sy)~--~(\sstartx-0.5,~\sy);~
 \draw[#9]~(\slengx+\sstartx+1.5-0.5,~\sy)~--~(\slengx+\sstartx+1.5+\sdxt-0.5,\sy);~
  }
\NewDocumentCommand \rcell { m m }
  {
%
~\pgfmathsetmacro{\x}{#1}~
 \pgfmathsetmacro{\y}{#2}~
 \draw (\x+1,\y)~node[aallele]{A};~
 \draw (\x+2,\y)~node[aallele]{C};~
 \draw (\x+3,\y)~node[aallele]{T};~
 \draw (\x+4,\y)~node[aallele]{C};~
 \gseg{\x}{4}{\y}{0.2}{0.5}{refseg}{1.5}{1}{reference};~
  }
\NewDocumentCommand \ncell { m m }
  {
%
~\pgfmathsetmacro{\x}{#1}~
 \pgfmathsetmacro{\y}{#2}~
 \draw[bncell](\x-2,\y+0.5)--(\x+7,\y+0.5)--~
 (\x+7,\y-1)--(\x-2,\y-1)--cycle;~
 \draw (\x+1,\y)~node[aallele]{A};~
 \draw (\x+2,\y)~node[ballele]{G};~
 \draw (\x+3,\y)~node[aallele]{T};~
 \draw (\x+4,\y)~node[aallele]{C};~
 \gseg{\x}{4}{\y}{0.2}{0.5}{mseg}{1.5}{1}{maternal};~
 \draw (\x+1,\y-0.5)~node[ballele]{T};~
 \draw (\x+2,\y-0.5)~node[aallele]{C};~
 \draw (\x+3,\y-0.5)~node[aallele]{T};~
 \draw (\x+4,\y-0.5)~node[ballele]{A};~
 \gseg{\x}{4}{\y-0.5}{0.2}{0.5}{pseg}{1.5}{1}{paternal};~
  }
\NewDocumentCommand \tcellone { m m }
  {
%
~\pgfmathsetmacro{\x}{#1}~
 \pgfmathsetmacro{\y}{#2}~
 \draw[btcell](\x-2,\y+0.5)--(\x+7,\y+0.5)--~
 (\x+7,\y-1.5)--(\x-2,\y-1.5)--cycle;~
 \draw (\x+1,\y)~node[aallele]{A};~
 \draw (\x+2,\y)~node[ballele]{G};~
 \draw (\x+3,\y)~node[aallele]{T};~
 \draw (\x+4,\y)~node[aallele]{C};~
 \gseg{\x}{4}{\y}{0.2}{0.5}{mseg}{1.5}{1}{maternal};~
 \draw (\x+1,\y-0.5)~node[aallele]{A};~
 \draw (\x+2,\y-0.5)~node[ballele]{G};~
 \draw (\x+3,\y-0.5)~node[aallele]{T};~
 \draw (\x+4,\y-0.5)~node[aallele]{C};~
 \gseg{\x}{4}{\y-0.5}{0.2}{0.5}{mseg}{1.5}{1}{maternal};~
 \draw (\x+1,\y-1)~node[ballele]{T};~
 \draw (\x+2,\y-1)~node[aallele]{C};~
 \draw (\x+3,\y-1)~node[aallele]{T};~
 \draw (\x+4,\y-1)~node[ballele]{A};~
 \gseg{\x}{4}{\y-1}{0.2}{0.5}{pseg}{1.5}{1}{paternal};~
  }
\NewDocumentCommand \tcellthree { m m }
  {
%
~\pgfmathsetmacro{\x}{#1}~
 \pgfmathsetmacro{\y}{#2}~
 \draw[btcell](\x-2,\y+0.5)--(\x+11,\y+0.5)--~
 (\x+11,\y-1)--(\x-2,\y-1)--cycle;~
 \draw (\x+1,\y)~node[aallele]{A};~
 \draw (\x+2,\y)~node[ballele]{G};~
 \gseg{\x}{2}{\y}{0.2}{0.5}{mseg}{1.5}{0}{maternal};~
 \gseg{\x+4}{0}{\y}{0.2}{0.5}{vseg}{0.5}{0.5}{variant};~
 \draw (\x+7,\y)~node[aallele]{T};~
 \draw (\x+8,\y)~node[aallele]{C};~
 \gseg{\x+6}{2}{\y}{0.2}{0.5}{mseg}{0}{1}{maternal};~
 \draw (\x+1,\y-0.5)~node[ballele]{T};~
 \draw (\x+2,\y-0.5)~node[aallele]{C};~
 \draw (\x+3,\y-0.5)~node[aallele]{T};~
 \draw (\x+4,\y-0.5)~node[ballele]{A};~
 \gseg{\x}{4}{\y-0.5}{0.2}{0.5}{pseg}{1.5}{1}{paternal};~
  }
\NewDocumentCommand \tcellfour { m m }
  {
%
~\pgfmathsetmacro{\x}{#1}~
 \pgfmathsetmacro{\y}{#2}~
 \draw[btcell](\x-2,\y+0.5)--(\x+15,\y+0.5)--~
 (\x+15,\y-1)--(\x-2,\y-1)--cycle;~
 \draw (\x+1,\y)~node[aallele]{A};~
 \draw (\x+2,\y)~node[ballele]{G};~
 \gseg{\x}{2}{\y}{0.2}{0.5}{mseg}{1.5}{0}{maternal};~
 \gseg{\x+4}{0}{\y}{0.2}{0.5}{vseg}{0.5}{0.5}{variant};~
 \draw (\x+7,\y)~node[aallele]{T};~
 \gseg{\x+6}{1}{\y}{0.2}{0.5}{mseg}{0}{0}{maternal};~
 \gseg{\x+9}{0}{\y}{0.2}{0.5}{vseg}{0.5}{0.5}{variant};~
 \draw (\x+12,\y)~node[aallele]{C};~
 \gseg{\x+11}{1}{\y}{0.2}{0.5}{mseg}{0}{1}{maternal};~
 \draw (\x+1,\y-0.5)~node[ballele]{T};~
 \draw (\x+2,\y-0.5)~node[aallele]{C};~
 \draw (\x+3,\y-0.5)~node[aallele]{T};~
 \draw (\x+4,\y-0.5)~node[ballele]{A};~
 \gseg{\x}{4}{\y-0.5}{0.2}{0.5}{pseg}{1.5}{1}{paternal};~
  }
\NewDocumentCommand \tcelltwo { m m }
  {
%
~\pgfmathsetmacro{\x}{#1}~
 \pgfmathsetmacro{\y}{#2}~
 \draw[btcell](\x-2,\y+0.5)--(\x+7,\y+0.5)--~
 (\x+7,\y-1.5)--(\x-2,\y-1.5)--cycle;~
 \draw (\x+1,\y)~node[aallele]{A};~
 \draw (\x+2,\y)~node[ballele]{G};~
 \draw (\x+3,\y)~node[aallele]{T};~
 \draw (\x+4,\y)~node[aallele]{C};~
 \gseg{\x}{4}{\y}{0.2}{0.5}{mseg}{1.5}{1}{maternal};~
 \draw (\x+1,\y-0.5)~node[aallele]{A};~
 \draw (\x+2,\y-0.5)~node[ballele]{G};~
 \draw (\x+3,\y-0.5)~node[aallele]{T};~
 \draw (\x+4,\y-0.5)~node[ballele]{G};~
 \gseg{\x}{4}{\y-0.5}{0.2}{0.5}{mseg}{1.5}{1}{maternal};~
 \draw (\x+1,\y-1)~node[ballele]{T};~
 \draw (\x+2,\y-1)~node[aallele]{C};~
 \draw (\x+3,\y-1)~node[aallele]{T};~
 \draw (\x+4,\y-1)~node[ballele]{A};~
 \gseg{\x}{4}{\y-1}{0.2}{0.5}{pseg}{1.5}{1}{paternal};~
  }
\NewDocumentCommand \tcellfive { m m }
  {
%
~\pgfmathsetmacro{\x}{#1}~
 \pgfmathsetmacro{\y}{#2}~
 \draw[btcell](\x-2,\y+0.5)--(\x+9.5,\y+0.5)--~
 (\x+9.5,\y-1.5)--(\x-2,\y-1.5)--cycle;~
 \draw (\x+1,\y)~node[aallele]{A};~
 \draw (\x+2,\y)~node[ballele]{G};~
 \draw (\x+3,\y)~node[aallele]{T};~
 \draw (\x+4,\y)~node[aallele]{C};~
 \gseg{\x}{4}{\y}{0.2}{0.5}{mseg}{1.5}{1}{maternal};~
 \draw (\x+1,\y-0.5)~node[aallele]{A};~
 \draw (\x+2,\y-0.5)~node[ballele]{G};~
 \draw (\x+3,\y-0.5)~node[aallele]{T};~
 \draw (\x+4,\y-0.5)~node[aallele]{C};~
 \gseg{\x}{4}{\y-0.5}{0.2}{0.5}{mseg}{1.5}{1}{maternal};~
 \draw (\x+1,\y-1)~node[ballele]{T};~
 \gseg{\x}{1}{\y-1}{0.2}{0.5}{pseg}{1.5}{0}{paternal};~
 \draw (\x+4.5,\y-1)~node[ballele]{A};~
 \draw (\x+5.5,\y-1)~node[aallele]{T};~
 \draw (\x+6.5,\y-1)~node[aallele]{C};~
 \gsegr{\x+3.5}{3}{\y-1}{0.2}{0.5}{pseg}{0.5}{1.5}{paternal};~
  }
%</thesis-floats>
%    \end{macrocode}
%
% 按照我工要求，页面中标题之下不少于一行。
% ^^A ======================================================================
% ^^A Module thesis-chapter: 章节标题格式（字号、缩进、间距、编号样式）（hit-thesis）
% ^^A ======================================================================
%    \begin{macrocode}
%<*thesis-chapter>
%    \end{macrocode}
%
% \changes{v3.2a}{0000/00/00}{标题字体、目录里章标题的字体与节号之后的间隔
%   迁到 expl3，四个重复的 \cs{ifboolexpr} 收敛成两个函数。}
% \cs{hit@title@font} 是标题字体，\cs{hit@chapter@entry@font} 是目录里章标题的字体
% （注意别跟 \file{book-toc} 里的 \cs{hit@chapter@entry@font} 弄混，那个只在
% \option{toc/sans-font} 打开时才定义）。\cs{hit@section@separator} 与
% \cs{hit@section@separator@shenzhen} 是节号与节名之间的间隔，两者只差哈尔滨本科算不算进
% \cs{quad}。四个都可展开，能直接写进 \cs{ctexset} 的取值。
%    \begin{macrocode}
\cs_new:Npn \hit@title@font
  { \bool_if:NTF \g__hit_arialtitle_bool { \sffamily } { \heiti } }
\cs_new:Npn \hit@chapter@entry@font
  { \bool_if:NTF \g__hit_arialtitle_bool { \sffamily \heiti } { \heiti } }
\cs_new:Npn \__hit_section_separator:n #1
  {
    \bool_lazy_or:nnTF
      {
        \bool_lazy_and_p:nn
          { \bool_if_p:N \g__hit_harbin_bool } { \bool_if_p:N \g__hit_master_bool }
      }
      { \bool_lazy_and_p:nn {#1} { \bool_if_p:N \g__hit_bachelor_bool } }
      { \quad } { \enspace }
  }
\cs_new:Npn \hit@section@separator@shenzhen
  { \__hit_section_separator:n { \bool_if_p:N \g__hit_shenzhen_bool } }
\cs_new:Npn \hit@section@separator
  {
    \__hit_section_separator:n
      {
        \bool_lazy_or_p:nn
          { \bool_if_p:N \g__hit_shenzhen_bool } { \bool_if_p:N \g__hit_harbin_bool }
      }
  }
%    \end{macrocode}
% \begin{macro}{\hit@chapter@title@format}
% 根据 \issue[220] 为英文摘要标题增加 chapterbold 的判断
% \changes{v3.1d}{2025/03/03}{根据 \issue[220] 为英文摘要标题增加 chapterbold 的判断}
%
% \changes{v3.2a}{0000/00/00}{章标题格式迁到 expl3。}
% 英文版全大写；中文版里英文摘要那一章不做悬挂缩进，其余按
% \option{style/chapter-hang} 决定。悬挂宽度必须用
% \cs{CTEX@chaptername}\cs{CTEX@chapter@aftername} 量，否则量到的是英文标题长度。
% 两个分支末尾的 \texttt{\string~} 是原写法里换行带出来的空格，改写时要保留。
%    \begin{macrocode}
\cs_new_protected:Npn \hit@chapter@title@format #1
  {
    \bool_if:NTF \g__hit_en_bool
      { \MakeUppercase {#1} ~ }
      {
        \hit@if@same@TF {#1} { \eabstractcname }
          { \bool_if:NT \g__hit_chapterbold_bool { \bfseries } #1 }
          {
            \bool_if:NT \g__hit_chapterhang_bool
              {
                \settowidth \hangindent { \CTEX@chaptername \CTEX@chapter@aftername }
                \int_set:Nn \hangafter { 1 }
              }
            #1
          } ~
      }
  }
%    \end{macrocode}
% \end{macro}
% 添加英文版目录chapter title控制
% \changes{v3.0.0}{2020/05/21}{添加英文版目录chapter title控制}
%    \begin{macrocode}
\NewDocumentCommand \hit@chapter@english@format { m }
  {
    \begin { center }
      \hit@title@font
      \legacy_if:nTF { @mainmatter } { \xiaoer [1.57481] } { \sanhao [1.57481] }
      \bfseries #1
    \end { center }
  }
\NewDocumentCommand \hit@chapter@english@name@format { m }
  { \begin { center } \MakeUppercase {#1}~\end { center } }
\NewDocumentCommand \hit@chapter@postdocformat { m }
  { \centering \hit@title@font \xiaoer [1.57481] \textbf {#1} }
%    \end{macrocode}
%
% 添加博后标题格式
% \changes{v3.0.0}{2020/05/21}{添加博后标题格式}
%
% \changes{v3.2a}{0000/00/00}{\cs{@afterheading} 的两个内核开关迁到 \cs{legacy\_if:nTF}}
%    \begin{macrocode}
\cs_set:Npn \@afterheading
  {
    \@nobreaktrue
    \everypar
      {
        \legacy_if:nTF { @nobreak }
          {
            \@nobreakfalse
            \clubpenalty 1
            \legacy_if:nF { @afterindent } { { \setbox \z@ \lastbox } }
          }
          { \clubpenalty 1 \everypar { } }
      }
  }
%    \end{macrocode}
%
% 设置一到四级标题、目录、书签格式。
%
% 按照规范中规定行距的mm数定义各级行距
% \changes{v3.0.10}{2020/10/28}{按照规范中规定行距的mm数定义各级行距}
% \changes{v3.1c}{2023/04/16}{按照威海本科模板调整章标题的上下间距。\issue[135]}
% \changes{v3.1d}{2025/03/03}{将判断换成可以展开的 \cs{ifboolexpe} 来解决 "Missing number" 的错误，\issue[200]}
% \changes{v3.1d}{2025/03/03}{发现 \issue[135] 中的间距改错位置了……}
% \changes{v3.1f}{0000/00/00}{威海本科模板章标题的上下间距与其他校区统一。}
%
% \changes{v3.2a}{0000/00/00}{三套章节格式各自取名，分派改用 expl3。}
% 章节格式分英文版、博后、其余三套。每套的 \cs{ctexset} 里有字面空格和控制空格，
% 按 \ref{sec:expl3-space} 的规矩留在 expl3 之外，只把分派放进去。
%    \begin{macrocode}
\cs_new_protected:Npn \hit@chapter@settings@english
  {
%
~\ctexset{%
 chapter={~
 afterindent=true,~
 pagestyle={hit@headings},~
 beforeskip={28.34646bp},%一个空行 1.57481 × 18
 afterskip={28.74646bp},%0.8应该不计算间距 0.8 × 18 + 0.57481×18
 aftername={\par},~
 format={\hit@chapter@english@format},~
 nameformat={\hit@chapter@english@name@format},%\center 会影响之后全局
 numberformat=\relax,~
 titleformat={\hit@chapter@title@format},~
 fixskip=true,~% 添加这一行去除默认间距
 %hang=true,
 },~
 section={~
 afterindent=true,~
 beforeskip={\hit@glue@skip{19.84252bp}{2.834646bp}},%上下空0.5行
 afterskip={\hit@glue@skip{19.84252bp}{2.834646bp}},~
 format={\hit@title@font\hit@glue@font@size{15bp}{21bp}{1.677267bp~\@minus 1.157391bp}\selectfont\bfseries},~
 aftername=\enspace,~
 fixskip=true,~
 break={},~
 },~
 subsection={~
 afterindent=true,~
 beforeskip={\hit@glue@skip{17.00787bp}{2.834646bp}},~
 afterskip={\hit@glue@skip{17.00787bp}{2.834646bp}},~
 format={\hit@title@font\hit@glue@font@size{14bp}{18bp}{1.842609bp~\@minus 0.9920497bp}\selectfont\bfseries},~
 aftername=\enspace,~
 fixskip=true,~
 break={},~
 },~
 subsubsection={~
 afterindent=true,~
 beforeskip={\hit@glue@skip{8.503937bp}{2.834646bp}},~
 afterskip={\hit@glue@skip{8.503937bp}{2.834646bp}},~
 format={\hit@title@font\normalsize},~
 aftername=\enspace,~
 fixskip=true,~
 break={},~
 },~
 paragraph/afterindent=true,~
 subparagraph/afterindent=true~
 }~
  }
\cs_new_protected:Npn \hit@chapter@settings@postdoc
  {
%
~\ctexset{%
 chapter={~
 afterindent=true,~
 pagestyle={hit@headings},~
 beforeskip={28.34646bp},%一个空行 1.57481 × 18
 afterskip={28.74646bp},%0.8应该不计算间距 0.8 × 18 + 0.57481×18
 aftername=\enspace,~
 format={\hit@chapter@postdocformat},%\center 会影响之后全局
 nameformat=\relax,~
 numberformat=\relax,~
 titleformat=\hit@chapter@title@format,~
 fixskip=true,~% 添加这一行去除默认间距
 %hang=true,
 },~
 section={~
 afterindent=true,~
 beforeskip={\hit@glue@skip{19.84252bp}{2.834646bp}},%上下空0.5行
 afterskip={\hit@glue@skip{19.84252bp}{2.834646bp}},~
 format={\hit@title@font\hit@glue@font@size{15bp}{21bp}{1.677267bp~\@minus 1.157391bp}\selectfont\bfseries},~
 aftername=\enspace,~
 fixskip=true,~
 break={},~
 },~
 subsection={~
 afterindent=true,~
 beforeskip={\hit@glue@skip{17.00787bp}{2.834646bp}},~
 afterskip={\hit@glue@skip{17.00787bp}{2.834646bp}},~
 format={\hit@title@font\hit@glue@font@size{14bp}{18bp}{1.842609bp~\@minus 0.9920497bp}\selectfont\bfseries},~
 aftername=\enspace,~
 fixskip=true,~
 break={},~
 },~
 subsubsection={~
 afterindent=true,~
 beforeskip={\hit@glue@skip{8.503937bp}{2.834646bp}},~
 afterskip={\hit@glue@skip{8.503937bp}{2.834646bp}},~
 format={\hit@title@font\normalsize},~
 aftername=\enspace,~
 fixskip=true,~
 break={},~
 },~
 paragraph/afterindent=true,~
 subparagraph/afterindent=true~
 }~
\par
  }
\cs_new_protected:Npn \hit@chapter@settings@other
  {
%
~\ctexset{%
 chapter={~
 afterindent=true,~
 pagestyle={hit@headings},~
 beforeskip={28.34646bp},%一个空行 1.57481 × 18
 afterskip={28.74646bp},%0.8应该不计算间距 0.8 × 18 + 0.57481×18
%    \end{macrocode}
%
% \changes{v3.1b}{2022/06/06}{根据部分规范来设置序号和标题之间的间距}
% 本部硕士，深圳本科的序号与标题间距为两个半角空格，其余设置为一个半角空格，如果其他规范中有两个半角空格，欢迎提 \issue
%    \begin{macrocode}
 aftername={\hit@section@separator@shenzhen},~
%    \end{macrocode}
%
% 添加本科生章标题加粗选项
% \changes{v3.0.01}{2020/06/06}{添加本科生章标题加粗选项}
%    \begin{macrocode}
 format={\centering\hit@title@font\xiaoer[1.57481]\hit@if@option@T{chapterbold}{\hit@if@option@T{bachelor}{\bfseries}}},%\center 会影响之后全局
 nameformat=\relax,~
 numberformat=\relax,~
 titleformat=\hit@chapter@title@format,~
 fixskip=true,~% 添加这一行去除默认间距
 %hang=true,
 },~
 section={~
 afterindent=true,~
 beforeskip={\hit@glue@skip{19.84252bp}{2.834646bp}},%上下空0.5行
 afterskip={\hit@glue@skip{19.84252bp}{2.834646bp}},~
 format={\hit@title@font\hit@glue@font@size{15bp}{21bp}{1.677267bp~\@minus 1.157391bp}\selectfont},~
%    \end{macrocode}
%
% \changes{v3.1b}{2022/06/06}{根据部分规范来设置序号和标题之间的间距}
%    \begin{macrocode}
 aftername={\hit@section@separator},~
 fixskip=true,~
 break={},~
 },~
 subsection={~
 afterindent=true,~
 beforeskip={\hit@glue@skip{17.00787bp}{2.834646bp}},~
 afterskip={\hit@glue@skip{17.00787bp}{2.834646bp}},~
 format={\hit@title@font\hit@glue@font@size{14bp}{18bp}{1.842609bp~\@minus 0.9920497bp}\selectfont},~
%    \end{macrocode}
%
% \changes{v3.1b}{2022/06/06}{根据部分规范来设置序号和标题之间的间距}
% \changes{v3.1d}{2025/03/03}{\issue[220] 说来就来，本部本科也改成两个空格了。但是章标题的间距没改。}
%    \begin{macrocode}
 aftername={\hit@section@separator},~
 fixskip=true,~
 break={},~
 },~
 subsubsection={~
 afterindent=true,~
 beforeskip={\hit@glue@skip{8.503937bp}{2.834646bp}},~
 afterskip={\hit@glue@skip{8.503937bp}{2.834646bp}},~
 format={\hit@title@font\normalsize},~
%    \end{macrocode}
%
% \changes{v3.1b}{2022/06/06}{根据部分规范来设置序号和标题之间的间距}
%    \begin{macrocode}
 aftername={\hit@section@separator},~
 fixskip=true,~
 break={},~
 },~
 paragraph/afterindent=true,~
 subparagraph/afterindent=true~
 }~
  }
\bool_if:NTF \g__hit_en_bool
  { \hit@chapter@settings@english }
  {
    \bool_if:NTF \g__hit_postdoc_bool
      { \hit@chapter@settings@postdoc } { \hit@chapter@settings@other }
  }
%    \end{macrocode}
%
% 此处添加英文版和博后的title设置
% \changes{v3.0.0}{2020/05/21}{此处添加英文版和博后的title设置}
% 设置附表、附录格式。
% \changes{v1.0.13}{2018/04/05}{此处添加中文目录中Abstract是否均为大写选项}
% \changes{v2.1.1}{2020/05/04}{添加chapter功能，使页眉、TOC不受protect的命令影响}
%
% \changes{v3.2a}{0000/00/00}{附录章标题的两套定义各自取名，分派改用 expl3。}
% \cs{hit@appendix@chapter} 分英文版与中文版两套，只有一套会真的执行。
%
% \cs{ExplSyntaxOn} 写在 \cs{hit@if@option@TF} 的花括号参数里不起作用：参数在读到
% \cs{hit@if@option@TF} 的时候就按外层的 catcode 记号化完了。所以两套定义各自取名放在
% 外面，再用 \cs{hit@if@option@TF} 选一个 \cs{cs\_new\_eq:NN} 过去。
%    \begin{macrocode}
\NewDocumentCommand \__hit_appendix_chapter_en:w { s o m o }
  {
    \IfBooleanT {#1}
      {
        \phantomsection
        \IfValueTF {#2} { \markboth {#2} {#2} } { \markboth {#3} {#3} }
        \par
        \hit@if@same@TF {#3} { \cabstractcname }
          { \relax }
          {
            \addcontentsline { toc } { chapter }
              { \texorpdfstring { \hit@chapter@entry@font #3 } {#3} }
          }
        \par
        \IfValueT {#4}
          {
            \addcontentsline { toe } { chapter }
              { \texorpdfstring { \bfseries #4 } {#4} }
          }
        \hit@chapter * {#3}
      }
  }
%    \end{macrocode}
%
% 添加英文版附表格式
% \changes{v3.0.0}{2020/05/21}{添加英文版附表格式}
%    \begin{macrocode}
\NewDocumentCommand \__hit_appendix_chapter_zh:w { s o m o }
  {
    \IfBooleanT {#1}
      {
        \phantomsection
        \IfValueTF {#2} { \markboth {#2} {#2} } { \markboth {#3} {#3} }
        \hit@if@same@TF {#3} { \eabstractcname }
          {
            \addcontentsline { toc } { chapter }
              {
                \texorpdfstring
                  { \hit@chapter@entry@font \hit@if@option@TF { absupper } { \MakeUppercase {#3} } {#3} }
                  {#3}
              }
          }
          {
            \addcontentsline { toc } { chapter }
              { \texorpdfstring { \hit@chapter@entry@font #3 } {#3} }
          }
        \IfValueT {#4}
          {
            \addcontentsline { toe } { chapter }
              { \texorpdfstring { \bfseries #4 } {#4} }
          }
        \hit@chapter * {#3}
      }
  }
\hit@if@option@TF { en }
  { \cs_new_eq:NN \hit@appendix@chapter \__hit_appendix_chapter_en:w }
  { \cs_new_eq:NN \hit@appendix@chapter \__hit_appendix_chapter_zh:w }
\NewDocumentCommand \BiAppChapter { m m }
  {
\phantomsection
 \chapter{#1}~
%    \end{macrocode}
%
% 此处添加保护选项
% \changes{v1.0.13}{2018/04/05}{添加\cs{texorpdfstring}命令去除书签中带有格式时的警告}
%    \begin{macrocode}
 \addcontentsline{toe}{chapter}{\texorpdfstring{\bfseries \xiaosi Appendix~\thechapter\nobreakspace{}\nobreakspace{}#2}{Appendix~\thechapter\nobreakspace{}\nobreakspace{}#2}}~
  }
%    \end{macrocode}
%
% 设置章节命令。s: 星号，表示在目录中出不出现序号。m: 必须要有的选项，中文章
% 节名称也即目录中名称，页眉中名称，书签中的名称。o: 可选内容，没有就默认是正
% 文章节，如果有，则是英文目录中显示的内容。
%    \begin{macro}{\chapter}
%    \begin{macrocode}
\cs_set_eq:NN \hit@chapter \chapter
\RenewDocumentCommand \chapter { s o m o }
  {
    \hit@clear@page \phantomsection
    \IfBooleanTF {#1}
      {% \chapter*
        \IfNoValueTF {#2}
          { \hit@chapter * {#3} }
          { \hit@chapter * [#2] {#3} }
        \IfValueT {#4}
          {
%    \end{macrocode}
%
% 此处添加保护选项
% \changes{v1.0.13}{2018/04/05}{添加\cs{texorpdfstring}命令去除书签中带有格式时的警告}
%    \begin{macrocode}
            \addcontentsline { toe } { chapter }
              { \texorpdfstring { \bfseries #4 } {#4} }
          }
      }
      {% \chapter
        \IfNoValueTF {#2}
          { \hit@chapter {#3} }
          { \hit@chapter [#2] {#3} }
        \IfValueT {#4}
          {
%    \end{macrocode}
%
% 此处需删除章节的空白
% \changes{v1.0.05}{2017/09/20}{添加\cs{ignorespaces}选项，矫正英文目录多出一个空白而无法对齐的bug}
% 此处添加保护选项
% \changes{v1.0.13}{2018/04/05}{添加\cs{texorpdfstring}命令去除书签中带有格式时的警告}
% \changes{v3.1b}{2022/06/06}{添加 \cs{hit@english@toc@chapter@name} 来修正英文目录中的附录显示。}
% \changes{v3.1d}{2025/03/03}{根据 \issue[182] 修改英文目录没有对齐的问题。}
% \changes{v3.1d}{2025/03/03}{根据 \issue[209] 修改英文目录有序号的章/附录前有空格的问题。}
% \changes{v3.2a}{0000/00/00}{章节命令迁到 expl3，顺带去掉英文目录章条目末尾多出的
%   一个空格，它来自源码里 \cs{addcontentsline} 参数中的换行}
%    \begin{macrocode}
            \addcontentsline { toe } { chapter }
              {
                \texorpdfstring
                  {
                    \bfseries \relax \hit@english@toc@chapter@name \thechapter
                    \hspace { 0.5em } \ignorespaces #4
                  }
                  { \hit@english@toc@chapter@name \thechapter \space \ignorespaces #4 }
              }
          }
      }
  }
%    \end{macrocode}
%    \end{macro}
%    \begin{macro}{\section}
%    \begin{macrocode}
\cs_set_eq:NN \hit@section \section
\RenewDocumentCommand \section { s o m o }
  {
    \IfBooleanTF {#1}
      {% \section*
        \IfNoValueTF {#2}
          { \hit@section * {#3} }
          { \hit@section * [#2] {#3} }
        \IfValueT {#4} { \addcontentsline { toe } { section } {#4} }
      }
      {% \section
        \IfNoValueTF {#2}
          { \hit@section {#3} }
          { \hit@section [#2] {#3} }
        \IfValueT {#4}
          {
%    \end{macrocode}
%
% 此处需删除章节的空白
% \changes{v1.0.05}{2017/09/18}{添加\cs{ignorespaces}选项，矫正英文目录多出一个空白而无法对齐的bug}
%    \begin{macrocode}
            \addcontentsline { toe } { section }
              { \protect \numberline { \use:c { thesection } } \ignorespaces #4 }
          }
      }
  }
%    \end{macrocode}
%    \end{macro}
%    \begin{macro}{\subsection}
%    \begin{macrocode}
\cs_set_eq:NN \hit@subsection \subsection
\RenewDocumentCommand \subsection { s o m o }
  {
    \IfBooleanTF {#1}
      {% \subsection*
        \IfNoValueTF {#2}
          { \hit@subsection * {#3} }
          { \hit@subsection * [#2] {#3} }
        \IfValueT {#4} { \addcontentsline { toe } { subsection } {#4} }
      }
      {% \subsection
        \IfNoValueTF {#2}
          { \hit@subsection {#3} }
          { \hit@subsection [#2] {#3} }
        \IfValueT {#4}
          {
%    \end{macrocode}
%
% 此处需删除章节的空白
% \changes{v1.0.05}{2017/09/18}{添加\cs{ignorespaces}选项，矫正英文目录多出一个空白而无法对齐的bug}
%    \begin{macrocode}
            \addcontentsline { toe } { subsection }
              { \protect \numberline { \use:c { thesubsection } } \ignorespaces #4 }
          }
      }
  }
%    \end{macrocode}
%    \end{macro}
%    \begin{macro}{\subsubsection}
%    \begin{macrocode}
\cs_set_eq:NN \hit@subsubsection \subsubsection
\RenewDocumentCommand \subsubsection { s o m o }
  {
    \IfBooleanTF {#1}
      {% \subsubsection*
        \IfNoValueTF {#2}
          { \hit@subsubsection * {#3} }
          { \hit@subsubsection * [#2] {#3} }
        \IfValueT {#4} { \addcontentsline { toe } { subsubsection } {#4} }
      }
      {% \subsubsection
        \IfNoValueTF {#2}
          { \hit@subsubsection {#3} }
          { \hit@subsubsection [#2] {#3} }
        \IfValueT {#4}
          {
%    \end{macrocode}
%
% 此处需删除章节的空白
% \changes{v1.0.05}{2017/09/18}{添加\cs{ignorespaces}选项，矫正英文目录多出一个空白而无法对齐的bug}
%    \begin{macrocode}
            \addcontentsline { toe } { subsubsection }
              { \protect \numberline { \use:c { thesubsubsection } } \ignorespaces #4 }
          }
      }
  }
%    \end{macrocode}
%    \end{macro}
%
% \subsubsection{定义封面}
% \label{sec:cov}
% 封面信息。
% \changes{v1.0.11}{2018/03/07}{更改的中文标题，根据反馈，在封面中标题需要自由换行且不能影响到原创性声明。此处额外设置了一个变量cover-title-zh。}
%    \begin{macrocode}
%</thesis-chapter>
%    \end{macrocode}
% ^^A ======================================================================
% ^^A Module thesis-meta: 元信息字段（\hitsetup 接受的 key）（hit-thesis）
% ^^A ======================================================================
%    \begin{macrocode}
%<*thesis-meta>
% \changes{v3.2a}{0000/00/00}{meta 模块抽取后，etoolbox 需提前加载以使用 \cs{AfterPreamble}}
\RequirePackage{etoolbox}
% \changes{v3.2a}{0000/00/00}{\cs{hitsetup} 的底层机制挪到 \file{src/shared/hit-meta.dtx}}

\hit@define@term{secret-level} %密级
\hit@define@term{classified-index-national}  %国内图书分类号
\hit@define@term{classified-index-international}  %国际图书分类号

\hit@define@title{zh}
\hit@define@term{discipline-zh} %中文学科
\hit@define@term{author-zh} %中文作者
\hit@define@term{supervisor-zh} %中文导师
\hit@define@term{associate-supervisor-zh} %中文副导师
\hit@define@term{co-supervisor-zh}%中文联合导师
\hit@define@term{affiliation-zh}%中文院系
\hit@define@term{specialty-zh}%中文专业
%    \end{macrocode}
%
% \changes{v3.1b}{2022/06/06}{增加深圳封面最下方的时间信息}
% \changes{v3.1d}{2025/03/03}{根据 \issue[226]，将 \texttt{szshortcdate} 改为 \texttt{firstpagecdate}，只要赋值就会显示。}
% 根据 \issue[95] 增加一个变量 \texttt{szshortcdate} 来存储答辩时间。
% 根据 \issue[226]，将 \texttt{szshortcdate} 改为 \texttt{firstpagecdate}，只要赋值就会显示。
%    \begin{macrocode}
\hit@define@term{cover-date-zh}
\hit@define@term{date-zh}
\hit@define@term{report-number}
\hit@define@term{station-zh}
%    \end{macrocode}
%
% 添加博后封皮选项
% \changes{v3.0.12}{2020/11/02}{添加博后封皮选项}
%    \begin{macrocode}
\hit@define@term{work-finish-date-zh}
\hit@define@term{report-submit-date-zh}
\hit@define@term{research-start-date-zh}
\hit@define@term{research-end-date-zh}
%    \end{macrocode}
%
% 添加博后封皮的title选项
% \changes{v3.0.0}{2020/05/21}{添加博后封皮的title选项}
%    \begin{macrocode}

\hit@define@term{student-id}%
\hit@define@term{student-type-zh}%



\hit@define@title{en}
\hit@define@term{discipline-en} %英文学科
\hit@define@term{author-en} %英文作者
\hit@define@term{supervisor-en} %英文导师
\hit@define@term{associate-supervisor-en} %英文副导师
\hit@define@term{co-supervisor-en} %英文联合导师
\hit@define@term{affiliation-en}
\hit@define@term{specialty-en}
\hit@define@term{date-en}
\hit@define@term{student-type-en}
\cs_new_protected:Npn \hit@collect@abstract@zh #1 { \cs_gset:Npn \hit@abstract@zh {#1} }
\cs_new_protected:Npn \hit@collect@abstract@en #1 { \cs_gset:Npn \hit@abstract@en {#1} }
\NewDocumentEnvironment { abstract-zh } { } { \Collect@Body \hit@collect@abstract@zh } { }
\NewDocumentEnvironment { abstract-en } { } { \Collect@Body \hit@collect@abstract@en } { }
\NewDocumentEnvironment { cabstract } { }
  {
    \hit@warn@renamed@environment { cabstract } { abstract-zh }
    \Collect@Body \hit@collect@abstract@zh
  }
  { }
\NewDocumentEnvironment { eabstract } { }
  {
    \hit@warn@renamed@environment { eabstract } { abstract-en }
    \Collect@Body \hit@collect@abstract@en
  }
  { }
%    \end{macrocode}
%
% \changes{v3.2a}{0000/00/00}{关键词的逗号拆分挪到 \file{src/shared/hit-meta.dtx}}
% 关键词的拆分机制见 \file{src/shared/hit-meta.dtx}。
%    \begin{macrocode}
\hit@parse@keywords{keywords-zh}{keywords@separator@zh}
\hit@parse@keywords{keywords-en}{keywords@separator@en}
\keys_define:nn { hit / class }
  {
    design .code:n =
      {
        \str_if_eq:nnTF {#1} { true }
          { \bool_gset_true:N  \g__hit_design_bool }
          { \bool_gset_false:N \g__hit_design_bool }
      } ,
    design .default:n = { true } ,
  }
%    \end{macrocode}
%
% \changes{v3.2a}{0000/00/00}{\cs{hitsetup} 改用 l3keys 做入口，加一层按职责
%   划分的子族路由。子族只做转发，键由 \texttt{hit/class} 族处理，行为
%   与改动前一致；后续再逐族把键迁过来。}
% \cs{hitsetup} 的入口换成 l3keys。根族只做路由，把 \option{info}、\option{style}
% 等分派到对应子族；子族一律把键原样转回 \texttt{hit/class}
% 族，所以两种写法结果相同：
% \begin{latex}
% \hitsetup{title-zh={...}}              % 扁平写法，一直支持
% \hitsetup{info={title-zh={...}}}       % 新的分族写法
% \end{latex}
% 根族的 \option{unknown} 兜底负责扁平写法，子族的兜底负责分族写法。等某个键真正
% 迁到 l3keys 之后，它就不再走兜底，两条路自然收敛。兜底本身、三态取值与布尔键的
% 声明机制都在 \file{src/shared/hit-meta.dtx}。
%
% \changes{v3.2a}{0000/00/00}{元信息字段正式注册进 \texttt{hit/info} 族，
%   不再依赖子族兜底转发。}
% 元信息字段名单。这批由 \cs{hit@define@term} 与 \cs{hit@parse@keywords} 批量生成，
% 每个字段同时有一个同名命令（\cs{title-zh} 等），键的处理器直接调用那个命令，
% 所以三种写法结果一致：\cs{hitsetup} 的扁平写法、分族写法、以及命令写法。
%    \begin{macrocode}
\hit@declare@info@field
  {
    secret-level , classified-index-national , classified-index-international , title-zh ,
    discipline-zh , author-zh , supervisor-zh , associate-supervisor-zh , co-supervisor-zh ,
    affiliation-zh , specialty-zh , cover-date-zh , date-zh , report-number , station-zh ,
    work-finish-date-zh , report-submit-date-zh , research-start-date-zh , research-end-date-zh
    , student-id , student-type-zh , title-en , discipline-en , author-en , supervisor-en ,
    associate-supervisor-en , co-supervisor-en , affiliation-en , specialty-en , date-en ,
    student-type-en , keywords-zh , keywords-en
  }
%    \end{macrocode}
%
% \changes{v3.2a}{0000/00/00}{\file{cfg} 里的可本地化字符串做成 \texttt{hit/const}
%   族的键，用户在导言区即可改写，不必去动生成物。}
% 可本地化字符串的名单。默认值在类文件末尾给出，\cs{hitsetup} 在导言区执行，
% 所以用户设的值覆盖默认值。键名怎么由宏名得来见 \file{src/shared/hit-meta.dtx}。
%    \begin{macrocode}
\hit@declare@constant
  {
    acknowledgement-heading-zh , acknowledgement-heading-en , authorization-heading-zh ,
    authorization-heading-en , authorization-text-zh , bachelor-authorization-text-zh ,
    authorization-usage-heading-zh , author-signature-field-label-zh ,
    bachelor-affiliation-field-label-zh , bachelor-date-field-label-zh ,
    bachelor-check-date-field-label-zh , bachelor-student-field-label-zh ,
    bachelor-document-type-zh , bachelor-degree-level-zh ,
    bachelor-teacher-comment-field-label-zh , bachelor-teacher-signature-field-label-zh ,
    bibliography-heading-en , affiliation-field-label-zh , associate-supervisor-field-label-zh ,
    author-field-label-zh , brace-left-zh , brace-right-zh , co-supervisor-field-label-zh ,
    date-field-label-zh , degree-of-discipline-format-zh , degree-field-label-zh ,
    keywords-separator-zh , keywords-field-label-zh , conclusion-heading-zh ,
    conclusion-heading-en , corresponding-address-heading-zh , school-zh , school-field-label-zh
    , specialty-field-label-zh , supervisor-field-label-zh , document-type-zh ,
    thesis-title-prefix-zh , degree-level-zh , date-blank-zh , declaration-heading-zh ,
    bachelor-declaration-heading-zh , declaration-text-zh , denotation-heading-zh ,
    denotation-heading-en , eqdenote-label-zh , eqdenote-label-en , eqdenote-dash ,
    doctor-publication-heading-zh , affiliation-field-label-en ,
    associate-supervisor-field-label-en , author-field-label-en , brace-left-en , brace-right-en
    , co-supervisor-field-label-en , date-field-label-en , degree-of-discipline-format-en ,
    degree-field-label-en , figure-prefix-en , table-prefix-en , keywords-separator-en ,
    school-en , school-field-label-en , specialty-field-label-en , supervisor-field-label-en ,
    degree-level-en , degree-level-attribute-en , date-field-label-front-zh ,
    graduate-affiliation-field-label-zh , graduate-affiliation-major-field-label-zh ,
    graduate-major-field-label-zh , graduate-date-field-label-zh ,
    graduate-enroll-date-field-label-zh , graduate-student-id-field-label-zh ,
    graduate-student-field-label-zh , graduate-supervisor-field-label-zh ,
    graduate-thesis-field-label-zh , harbin-bachelor-school-bottom-mark ,
    harbin-school-bottom-mark , harbin-campus , hi , index-heading-en ,
    classified-index-international-field-label-zh , classified-index-national-field-label-zh ,
    postdoc-udc-field-label , postdoc-author-field-label-zh ,
    postdoc-classified-index-field-label-zh , postdoc-report-spaced-zh , postdoc-report-zh ,
    postdoc-research-end-date-field-label-zh , postdoc-work-finish-date-field-label-zh ,
    postdoc-specialty-field-label-zh , postdoc-number-field-label-zh ,
    postdoc-station-field-label-zh , postdoc-secret-level-field-label-zh , postdoc-separator ,
    postdoc-research-start-date-field-label-zh , postdoc-report-submit-date-field-label-zh ,
    postdoc-co-supervisor-field-label-zh , publication-heading-zh , publication-heading-en ,
    resolution-heading-zh , resolution-heading-en , resume-heading-zh , resume-heading-en ,
    school-id , school-id-field-label-zh , bachelor-second-authorization-text-zh ,
    secret-level-field-label-zh , shenzhen-doctor-midterm-note , shenzhen-school-bottom-mark ,
    shenzhen-campus , shenzhen-campus-postdoc , bachelor-affiliation-field-label-shenzhen-zh ,
    bachelor-affiliation-field-label-weihai-zh , bachelor-major-field-label-zh ,
    bachelor-student-id-field-label-zh , bachelor-supervisor-field-label-zh ,
    bachelor-thesis-field-label-zh , inline-separator , stage-document-type , stage-midterm ,
    stage-opening , supervisor-signature-field-label-zh , title-separator-zh ,
    title-separator-en , weihai-campus
  }
%    \end{macrocode}
%
% \changes{v3.2a}{0000/00/00}{参考文献的五个选项迁入 \texttt{hit/bib} 族，
%   改用连字符命名。旧名转回 \texttt{hit/class} 族处理，行为不变。}
% \changes{v3.2a}{0000/00/00}{目录的四个选项迁入 \texttt{hit/toc} 族，改用连字符
%   命名，其中三个改为三态取值。}
% \changes{v3.2a}{0000/00/00}{排版开关迁入 \texttt{hit/style} 族并改用连字符命名。}
% 三个族的键。排版开关都是普通布尔，用 \cs{\_\_hit\_style\_pair:nn} 成对声明，
% 新名设的是旧标志位本身，所以两套名字天然一致。目录的前三个键用三态取值，
% \texttt{auto} 的含义分别是：\option{english} 博士排、其余不排；\option{blank}
% 英文版不插、其余插；\option{four} 一律不开。参考文献的五个键此前分散在
% \file{book} 与 \file{art} 两个模块、命名也不统一（\texttt{citetwo} 看不出是
% 控制什么的），新旧两套名字设的是同一批内部变量，所以可以混用。
%    \begin{macrocode}
\keys_define:nn { hit / style }
  {
    ziju .choices:nn = { regu , word }
      {
        \bool_gset_false:N \g__hit_ziju_regu_bool
        \bool_gset_false:N \g__hit_ziju_word_bool
        \bool_gset_true:c { g__hit_ziju_ \l_keys_choice_tl _bool }
      } ,
    ziju .value_required:n = true ,
    headrule .code:n = { \__hit_tristate_set:Nn \g__hit_head_rule_tl {#1} } ,
    pagenumberline .meta:nn = { hit / style / pagenumberline } {#1} ,
  }
\keys_define:nn { hit / style / pagenumberline }
  {
    mainmatter  .code:n = { \__hit_tristate_set:Nn \g__hit_page_number_line_main_tl  {#1} } ,
    frontmatter .code:n = { \__hit_tristate_set:Nn \g__hit_page_number_line_front_tl {#1} } ,
  }
\clist_map_inline:nn
  {
    { chapter-bold } { chapterbold } ,
    { chapter-hang } { chapterhang } ,
    { sans-title } { arialtitle } ,
    { caption-center-last } { capcenterlast } ,
    { subcaption-center-last } { subcapcenterlast } ,
    { abstract-upper } { absupper } ,
    { etitle-auto } { etitleauto } ,
    { glue } { glue } ,
    { open-right } { openright } ,
    { library } { library } ,
    { ragged-bottom } { raggedbottom } ,
  }
  { \__hit_style_pair:nn #1 }
%    \end{macrocode}
%
%    \begin{macrocode}
\keys_define:nn { hit / toc }
  {
    english .code:n = { \__hit_tristate_set:Nn \g__hit_toc_english_tl {#1} } ,
    blank   .code:n = { \__hit_tristate_set:Nn \g__hit_toc_blank_tl   {#1} } ,
    four    .code:n = { \__hit_tristate_set:Nn \g__hit_toc_four_tl    {#1} } ,
    sans-font .choice: ,
    sans-font / true  .code:n = { \bool_gset_true:N \g__hit_arialtoc_bool } ,
    sans-font / false .code:n = { \bool_gset_false:N \g__hit_arialtoc_bool } ,
    sans-font .default:n = true ,
  }
%    \end{macrocode}
%
%    \begin{macrocode}
\keys_define:nn { hit / bib }
  {
    adjacent-cite .choices:nn =
      { comma , endash }
      {
        \bool_gset_false:N \g__hit_citetwo_comma_bool
        \bool_gset_false:N \g__hit_citetwo_endash_bool
        \bool_gset_true:c { g__hit_citetwo_ \l_keys_choice_tl _bool }
      } ,
    adjacent-cite .value_required:n = true ,
    label-align .code:n = { \str_set:Nn \g_bib_labelalign_str {#1} } ,
    author-case .code:n = { \str_if_eq:nnT {#1} { full }
      { \bool_set_true:N \g_bib_authorupper_bool } } ,
    max-author .code:n =
      {
        \int_set:Nn \g_bib_maxauthor_int {#1}
        \int_if_zero:nT { \g_bib_maxauthor_int } { \int_gdecr:N \g_bib_maxauthor_int }
      } ,
    split-item .choice: ,
    split-item / true  .code:n = { \bool_gset_true:N \g__hit_splitbibitem_bool } ,
    split-item / false .code:n = { \bool_gset_false:N \g__hit_splitbibitem_bool } ,
    split-item .default:n = true ,
  }
%    \end{macrocode}
%
% \changes{v3.2a}{0000/00/00}{加旧名到新名的映射表与废弃警告。旧名照常生效，
%   只是编译时提示新写法。}
% 兼容期的处理。旧名一律照常生效，只在日志里给一条提示。三类提示分开：
% 改过名的选项逐个报，元信息字段的扁平写法合并成一条报，
% \cs{documentclass} 里传排版类选项单独报。
%    \begin{macrocode}
\prop_gset_from_keyval:Nn \g__hit_renamed_option_prop
  {
    arialtoc         = { toc = { sans-font = ... } } ,
    tocfour          = { toc = { four = ... } } ,
    tocblank         = { toc = { blank = ... } } ,
    engtoc           = { toc = { english = ... } } ,
    arialtitle       = { style = { sans-title = ... } } ,
    chapterbold      = { style = { chapter-bold = ... } } ,
    chapterhang      = { style = { chapter-hang = ... } } ,
    capcenterlast    = { style = { caption-center-last = ... } } ,
    subcapcenterlast = { style = { subcaption-center-last = ... } } ,
    absupper         = { style = { abstract-upper = ... } } ,
    etitleauto       = { style = { etitle-auto = ... } } ,
    glue             = { style = { glue = ... } } ,
    openright        = { style = { open-right = ... } } ,
    library          = { style = { library = ... } } ,
    raggedbottom     = { style = { ragged-bottom = ... } } ,
    zijv             = { style = { ziju = ... } } ,
    bsheadrule       = { style = { headrule = ... } } ,
    bsmainpagenumberline  = { style = { pagenumberline = { mainmatter = ... } } } ,
    bsfrontpagenumberline = { style = { pagenumberline = { frontmatter = ... } } } ,
    citetwo          = { bib = { adjacent-cite = ... } } ,
    splitbibitem     = { bib = { split-item = ... } } ,
    biblabelalign    = { bib = { label-align = ... } } ,
    bibauthorupper   = { bib = { author-case = ... } } ,
    bibmaxauthor     = { bib = { max-author = ... } } ,
  }
%    \end{macrocode}
%
% \changes{v3.2a}{0000/00/00}{元信息字段的旧名走废弃期：仍然可用，用一次提示一次}
% 旧名到新名的对照。旧名在 v3.1e 就是公开接口，写在用户的 \file{cover.tex} 里，
% 所以保留到 v4.1。
%    \begin{macrocode}
\hit@define@legacy@terms
  {
    statesecrets=secret-level ,
    natclassifiedindex=classified-index-national ,
    intclassifiedindex=classified-index-international ,
    ctitle=title-zh ,
    cxueke=discipline-zh ,
    cauthor=author-zh ,
    csupervisor=supervisor-zh ,
    cassosupervisor=associate-supervisor-zh ,
    ccosupervisor=co-supervisor-zh ,
    caffil=affiliation-zh ,
    csubject=specialty-zh ,
    firstpagecdate=cover-date-zh ,
    cdate=date-zh ,
    cnumber=report-number ,
    cpositionname=station-zh ,
    cfinishdate=work-finish-date-zh ,
    csubmitdate=report-submit-date-zh ,
    cstartdate=research-start-date-zh ,
    cenddate=research-end-date-zh ,
    cstudentid=student-id ,
    cstudenttype=student-type-zh ,
    etitle=title-en ,
    exueke=discipline-en ,
    eauthor=author-en ,
    esupervisor=supervisor-en ,
    eassosupervisor=associate-supervisor-en ,
    ecosupervisor=co-supervisor-en ,
    eaffil=affiliation-en ,
    esubject=specialty-en ,
    edate=date-en ,
    estudenttype=student-type-en ,
    ckeywords=keywords-zh ,
    ekeywords=keywords-en
  }
\hit@define@legacy@title{ctitlecover}{zh}{first}
\hit@define@legacy@title{ctitleone}{zh}{first}
\hit@define@legacy@title{ctitletwo}{zh}{second}
\hit@define@legacy@title{csubtitle}{zh}{subtitle}
\hit@define@legacy@title{esubtitle}{en}{subtitle}
\prop_gset_from_keyval:Nn \g__hit_renamed_load_option_prop
  {
    font     = { font-en = ... } ,
    cjk-font = { font-zh = ... } ,
  }
%    \end{macrocode}
%
% \changes{v3.2a}{0000/00/00}{提示文案、只报一次的机制与类选项扫描挪到
%   \file{src/shared/hit-meta.dtx}}
% 单双面由学位与校区决定，所以扫描类选项时额外查一下 \option{oneside} 与
% \option{twoside}。另外两个类没有这条。
%    \begin{macrocode}
\cs_set_protected:Npn \__hit_class_option_extra:
  {
    \clist_if_in:nVT { oneside , twoside } \l__hit_warn_message_tl
      {
        \exp_args:NV \__hit_warn_once:nnn \l__hit_warn_message_tl
          { manual-sides } { }
      }
  }
\__hit_check_class_options:
\keys_define:nn { hit }
  {
    info  .meta:nn = { hit / info  } {#1} ,
    style .meta:nn = { hit / style } {#1} ,
    toc   .meta:nn = { hit / toc   } {#1} ,
    bib   .meta:nn = { hit / bib   } {#1} ,
    const .meta:nn = { hit / const } {#1} ,
    theorem .code:n = { \hit@theorem@style {#1} } ,
    theorem-list .code:n = { \hit@theorem@list {#1} } ,
    theorem-star .code:n = { \hit@theorem@declare@star@list {#1} } ,
    unknown .code:n = { \__hit_legacy_key:n {#1} } ,
  }
\RenewDocumentCommand \hitsetup { +m } { \keys_set:nn { hit } {#1} }
%    \end{macrocode}
%
% \changes{v3.1d}{2025/03/03}{将 \cs{firstpagecdate} 的设置藏起来}
%    \begin{macrocode}
\AfterPreamble{%
  \ifdefempty{\hit@cover@date@zh}{\use:c{cover-date-zh}{\hit@date@zh}}{}%
}
%</thesis-meta>
%    \end{macrocode}
%
% \changes{v3.2a}{0000/00/00}{参考文献相关代码从 bookcls 拆出至 book-bib 模块}
% ^^A ======================================================================
% ^^A Module thesis-bib: 参考文献加载与 natbib/gbt7714 设置（hit-thesis）
% ^^A ======================================================================
%    \begin{macrocode}
%<*thesis-bib>
\RequirePackage[sort&compress]{gbt7714}
\RequirePackage{etoolbox}
\RequirePackage{xparse}
\cs_new:Npn \bibstyle@numerical
  { \bibpunct {[} {]} {,} {s} {,} { \textsuperscript {,} } }
%    \end{macrocode}
%
% 删除inline等定义，此命令与 \pkg{gbt7714} 冲突
% \changes{v3.0.15}{2021/05/06}{删除inline等定义，此命令与 \pkg{gbt7714} 冲突}
%    \begin{macrocode}
% \newcommand\bibstyle@authoryear{\bibpunct{(}{)}{;}{a}{,}{,}}
%    \end{macrocode}
%
% \changes{v2.0.06}{2018/12/05}{在\cs{inlinecite}内添加空格}
%    \begin{macrocode}
% \newcommand\bibstyle@inline{\bibpunct{[}{]}{,}{n}{,}{\hit@inline@separator}}
\citestyle{numerical}
\DeclareRobustCommand\inlinecite{\@inlinecite}
\cs_set:Npn \@inlinecite #1
  {
    \group_begin:
      \citestyle { inline }
      \cs_set_eq:NN \@cite \NAT@citenum
      \citep {#1}
    \group_end:
  }
\cs_set_eq:NN \onlinecite \inlinecite
%    \end{macrocode}
%
% \changes{v3.1a}{2022/05/17}{按照 \hitszthesis{} 生成深圳校区的参考文献样式 \texttt{hitszthesis.bst}。\issue[97]}
%    \begin{macro}{\@biblabel}
% 设置参考文献列表的序号中的括号为中文括号。
% \changes{v3.1b}{2022/06/06}{将 ［］前加上 \cs{mbox\{\}} 来获得正确的序号长度。} 
% 在［\#1］前加上 \cs{mbox\{\}} 来获得正确的序号长度，否则会因为 \XeLaTeX 下的标点挤压而计算错误，这里应该整个包裹起来，否则换行后缩进会有问题。
% \changes{v3.1d}{2025/03/03}{按着 \issue[211] 来修改，将 \cs{mbox\{\}} 包住整个 label。}
% \changes{v3.1d}{2025/03/03}{按照 \issue[215] 将哈尔滨硕博毕业论文的参考文献序号也改为全角。}
%    \begin{macrocode}
\bool_lazy_or:nnT
  { \bool_if_p:N \g__hit_shenzhen_bool }
  {
    \bool_lazy_and_p:nn
      { \bool_if_p:N \g__hit_harbin_bool }
      {
        \bool_lazy_or_p:nn
          { \bool_if_p:N \g__hit_master_bool } { \bool_if_p:N \g__hit_doctor_bool }
      }
  }
  { \cs_set:Npn \@biblabel #1 { \mbox { ［#1］ } } }
%    \end{macrocode}
% \end{macro}
% 
% \begin{environment}{thebibliography}
% 修改默认的 \env{thebibliography} 环境来适应间距的要求。
% \changes{v3.2a}{0000/00/00}{设置参考文献列表标号的对齐方式}
%    \begin{macrocode}
\RenewDocumentEnvironment{thebibliography}{m}{%
%    \end{macrocode}
%
% \changes{v3.1d}{2025/03/03}{更改完全角符号后的半角点号可能出现在行首，规避一下。}
% 这里我直接贴一个 issue 的链接：\href{https://github.com/CTeX-org/ctex-kit/issues/701}{CTeX-Kit \#701}
% \changes{v3.2a}{0000/00/00}{FullRight引入了太多bug, 取消之 (bst文件中加入了nolinebreak)}
%    \begin{macrocode}
  % \xeCJKDeclareCharClass{FullRight}{"002E}%
  \hit@clear@page\phantomsection%
  \hit@appendix@chapter*{\bibname}[\hit@bibliography@heading@en]
  \normalsize
  \list{\@biblabel{\@arabic\c@enumiv}}{%
    \str_if_eq:VnT \g_bib_labelalign_str {left} {
      \RenewDocumentCommand\makelabel{m}{##1\hfill}
    }
    \settowidth{\labelwidth}{\@biblabel{#1}}
%    \end{macrocode}
% \changes{v3.1b}{2022/06/06}{修改深圳校区的 \cs{labelsep} 为 0pt。}
%    \begin{macrocode}
    \hit@if@option@TF{shenzhen}{\dim_set:Nn \labelsep {0pt}}{\dim_set:Nn \labelsep {0.5em}}
    \dim_set:Nn \itemindent {0pt}
    \dim_set:Nn \leftmargin {\labelsep+\labelwidth}
    \dim_add:Nn \itemsep {-0.8em}
    \usecounter{enumiv}%
    \cs_set_eq:NN \p@enumiv \@empty
  \cs_set:Npn \theenumiv {\@arabic\c@enumiv}}%
  \sloppy\frenchspacing
%    \end{macrocode}
% \changes{v2.0.07}{2019/03/02}{添加flushbottom到thebibliography环境中}
%    \begin{macrocode}
  \flushbottom
%    \end{macrocode}
% \changes{v2.0.03}{2018/10/08}{添加参考文献分割开关}
%    \begin{macrocode}
  \hit@if@option@TF{splitbibitem}%
  {%
  \clubpenalty0
  \@clubpenalty \clubpenalty
  \widowpenalty0%
  \interlinepenalty-50%
  }%
  {%
  \clubpenalty4000
  \@clubpenalty \clubpenalty
  \widowpenalty4000%
  \interlinepenalty4000%
  }%
\sfcode`\.\@m}
{\cs_set_nopar:Npn \@noitemerr
  {\@latex@warning{Empty `thebibliography' environment}}%
\endlist\frenchspacing}
%    \end{macrocode}
% \end{environment}
%    \begin{macrocode}
\cs_set:Npn \NAT@cite #1#2#3
      {\ifNAT@swa\NAT@@open\if*#2*\else#2\NAT@spacechar\fi
        #1\NAT@@close\if*#3*\else\textsuperscript{#3}\fi\else#1\fi\endgroup}
%    \end{macrocode}
% \changes{v3.2a}{0000/00/00}{控制参考文献列表中作者是否全大写和最大显示数量}
%    \begin{macrocode}
%% 用以控制参考文献列表中单个作者的格式
\exp_args:NNe \cs_new:Nn \g_bib_author_cs:n {
    \bool_if:NTF \g_bib_authorupper_bool {
        \exp_not:N \text_uppercase:n {#1}
    } {
        {#1}
    }
}
%% 用以输出参考文献列表中的作者
\NewDocumentCommand{\bibauthor}{m}{
  \int_zero:N \l_tmpa_int
  \exp_args:Ne \tl_map_inline:nn {
    \exp_args:Ne \tl_tail:n {#1}
  } {
    \int_compare:nNnT
      \l_tmpa_int = \g_bib_maxauthor_int {
      \prg_break:n {\tl_head:N {#1}}
    }
    \int_incr:N \l_tmpa_int
    \g_bib_author_cs:n {##1}
  }
  \prg_break_point:
}
%    \end{macrocode}
%
% \changes{v3.2a}{0000/00/00}{\cs{citeup}（上标引用宏）从 book-floats 迁入 book-bib，归位至引用语义所属模块}
% 上标引用宏 \cs{citeup}（语义属引用,非浮动体）。
%    \begin{macrocode}
\NewDocumentCommand \citeup { m } { \textsuperscript { \cite {#1} } }
%</thesis-bib>
%    \end{macrocode}
